% =====================================================================
%  Constans [Hans Brøchner], „Det Kristelige og det Humane som
%  Modsætninger i vor Tids Bevidsthed“ I--III.
%  Nyt dansk Maanedsskrift, udg. af Vilhelm Møller, 2. Bind (April --
%  September 1871): I. pp. 289--299; II. pp. 433--456; III. pp. 492--510.
%  TRANSCRIPTION (Danish) -- diplomatic, from the page images.
%
%  ---- SOURCE -----------------------------------------------------------
%  Internet Archive item nyt-dansk-maanedsskrift-bind-2 (whole volume,
%  Nyt_dansk_Maanedsskrift_bind2.pdf, 581 PDF pp., 327118779 bytes, sha256
%  449b659b6c4b4fac3d76ee8835c55f0023536cfb83e6123988ead21b9870cdcb,
%  ABBYY FineReader text layer), downloaded 2026-09-27:
%  https://archive.org/details/nyt-dansk-maanedsskrift-bind-2
%  The installments are volume PDF pp. 295--305, 439--462, 498--516.
%  Kept as three extracts (qpdf, text layer kept) in ~/bibliotek/Brøchner, Hans/:
%    1871-kristelige-humane-I-nyt-dansk-maanedsskrift-2.pdf   sha256 cce88142...
%    1871-kristelige-humane-II-nyt-dansk-maanedsskrift-2.pdf  sha256 cfe72857...
%    1871-kristelige-humane-III-nyt-dansk-maanedsskrift-2.pdf sha256 0ad82bce...
%  Each page one 300 ppi RGB JPEG.  Fraktur.
%
%  ---- PAGE MAP ---------------------------------------------------------
%  Volume PDF = printed + 6 throughout; verified by eye on the folios.
%
%  ---- CONVENTIONS ------------------------------------------------------
%  Diplomatic; 1871 orthography as printed.  Letterspacing -> \emph;
%  antiqua against the Fraktur -> \textit.  Fraktur I/J resolved by the
%  word.  Quotation marks „...“ as printed.  Footnotes *) at the mark.
%  Printer's errors transcribed as printed and logged in a % comment at
%  the site.
%
%  ---- LAYOUT -----------------------------------------------------------
%  Each installment opens at a fresh page (upper part blank) under the
%  title in display blackletter over two lines, a short rule, and the
%  centred numeral I. / II. / III.; each closes with „Constans.“ set right
%  in bold Fraktur, then a centred double rule (I, III) or a single rule
%  (II).  Short centred rules divide the text within each installment.
%  One footnote (p. 436).  Antiqua: »pagani« (p. 294, in roman
%  guillemets), imprimatur (p. 440), Juste-milieu (p. 442).  Emphasis
%  (letterspacing, measured within the line): 119 runs, sparse in I,
%  dense in III; compounds spaced only in part are marked as printed.
%
%  ---- ERRATUM ----------------------------------------------------------
%  The volume's „Tilføielser og Rettelser“ (p. 572) corrects p. 444, line
%  4 from the foot: „Viden“ for „Veien“.  Transcribed as printed; logged
%  at the site.  (This copy has a reader's ink note to the same effect.)
%
%  Printer's errors logged at the site: p. 289 „for-|laret“; p. 292 a
%  broken e in „Synden“; p. 433 „Det er det;“; p. 443 „Stihed“; p. 449
%  „sandelige“; p. 451 no point after „(V. 16 f.)“; p. 452
%  „Anthropomofismer“ (4x); p. 500 „Anthropomosisme“; p. 506 „Saa|ledes“
%  broken without a hyphen; p. 507 no point after „(Calvin)“; p. 510 „religtøse“.
%
%  ---- ACCURACY  [2026-09-27] ------------------------------------------
%  Word-accuracy: every batch diffed against the embedded ABBYY layer
%  before splicing (ocrdiff.py --frag; offsets -288, -421, -456 in the
%  54-page working extract): 5 batches (pp. 289-299, 433-444, 445-456,
%  492-501, 502-510), 117 candidates, 0 transcription errors corrected,
%  117 the witness's fault, 0 unresolved.  See TRANSCRIPTION-PLAYBOOK.md
%  §3 step 4.  The transcribers took their words from that layer
%  corrected at the image, so the diff cannot see a layer error they
%  copied; the second reading exists for that class.
%  SECOND READING of all 54 pages, word by word and mark by mark at the
%  image (300 dpi; 600-900 dpi for any doubt), by five readers who had
%  transcribed nothing and were forbidden the text layer: 5 discrepancies,
%  all applied.  Words: p. 289 „saae“ -> „faae“ (a crossed f, not a long s:
%  the layer's f/s confusion, copied; confirmed at 600 dpi in the caller);
%  p. 510 „religiøse“ -> „religtøse“, a printer's error (t for i) that
%  had been silently normalised (confirmed at 600 dpi in the caller; now
%  logged).  Other: p. 494 „Tro“ letterspaced like „Viden“; p. 497 opens a
%  new paragraph (indent; p. 496 ends on a short line); p. 501 the
%  letterspacing runs into „in-“ of „indskrænkningen“.  The readers
%  settled the open points: p. 435 „Individets“ (a broken e, word
%  stands), p. 437 a new paragraph (indented), p. 499 „Karakteer“ (faint
%  but present); and found no sign of stale renders on pp. 445-456 or
%  502-510.  Result: 1 transcription word error and 1 silently corrected
%  printer's error in 54 pages, both corrected.
% =====================================================================
\documentclass[12pt,a4paper,openany]{book}
\usepackage[utf8]{inputenc}
\usepackage[T1]{fontenc}
\usepackage{textalpha}
\usepackage{libertinus}
\usepackage{libertinust1math}
\usepackage{graphicx}    % for \reflectbox (the old Danish ɔ: id-est mark)
\DeclareUnicodeCharacter{0254}{\reflectbox{c}}  % ɔ (U+0254) = reversed c
\usepackage[danish]{babel}
\usepackage[protrusion=true,expansion=true]{microtype}
\usepackage{setspace}
\usepackage[margin=1.2in]{geometry}
\usepackage{fancyhdr}
\setlength{\headheight}{28pt}
\addtolength{\topmargin}{-13.5pt}

% Footnotes as the 1877 printing sets them: *).
\renewcommand{\thefootnote}{*)}

% The article's numbered sections: the number alone, centred on its line.
\newcommand{\secnum}[1]{\par\begin{center}#1\end{center}\par\nobreak}

\usepackage{hyperref}
\hypersetup{
  pdftitle={Det Kristelige og det Humane som Modsætninger i vor Tids Bevidsthed},
  pdfauthor={Hans Brøchner},
  hidelinks
}

\pagestyle{fancy}
\fancyhf{}
\fancyhead[LE,RO]{\thepage}
\fancyhead[RE]{\textit{Constans [H. Brøchner]: Det Kristelige og det Humane}}
\fancyhead[LO]{\textit{Nyt dansk Maanedsskrift, 2.\ Bind}}

\setstretch{1.3}

\title{\textbf{Det Kristelige og det Humane\\ som Modsætninger i vor Tids Bevidsthed}\\[6pt]\large I--III}
\author{Constans [Hans Brøchner]}
\date{Nyt dansk Maanedsskrift, udg.\ af Vilhelm Møller,
  2.\ Bind (April--September 1871), pp.~289--299, 433--456, 492--510}

\usepackage{marginnote}
\newcommand{\opage}[1]{\marginnote{\footnotesize\textit{[#1]}}}
\newcommand{\apage}[1]{\marginnote{\footnotesize\textit{[#1]}}}

\begin{document}
\maketitle
\thispagestyle{empty}
\clearpage

% ===== I. (Nyt dansk Maanedsskrift 2, pp. 289--299) =====
% --- p. 289 ---
% p. 289: the upper part of the page is blank (the installment opens a new page; no preceding article, no running head); the head is set in a display blackletter (textura) in two lines, „Det Kristelige og det Humane som Modsætninger“ / „i vor Tids Bevidsthed.“; a short rule below it; then the numeral „I.“ in roman type, centred
\opage{289}\begin{center}\textbf{\large Det Kristelige og det Humane som Modsætninger i vor Tids Bevidsthed.}\end{center}

\secnum{I.}

Hvad der i vor Tid træder frem som en Kamp mellem
Tro og Viden, er, trods denne Modsætnings centrale Betydning,
dog kun et særligt Udtryk for en mere omfattende
Kamp: Kampen mellem det Kristelige og det Humane, og i
denne repræsenterer atter Kristendommen, som de positive
Religioners høieste og afsluttende Form, som deres reneste
og meest aandeliggjorte Udprægning, et mere Almeent: det
positivt Religiøse i dets Modsætning til det frie Menneskelige.
I den sidste Skikkelse har Modsætningen allerede i
Aartusinder været udpræget for Bevidstheden: den er brudt
frem, saasnart den menneskelige Tanke, dristig og stærk, havde
fundet sin Lov og sin Norm i Tanken selv, saasnart den i
Tilværelsen, som den nu selv vilde erkjende, ikke faae
% printer's error: „for-|laret“ across the line break (for „forklaret“; no k at the start of the next line); transcribed as printed
forlaret gjennem Mythens Gaader, havde fundet den samme
Fornuft, der er dens eget Væsen.

Allerede i Grækenlands Oldtid traadte den frie Tanke i
Modsætning til den Form af det positivt Religiøse, der var
given i den græske Folkereligion. De mange Guder med
deres Strid og Kamp, med deres menneskelige Lidenskaber,
% --- p. 290 ---
\opage{290}med deres menneskelige Skjæbner, kort sagt, med al deres
Endelighed, bleve et Anstød for den Tanke, der i den hele
Tilværelse kun søgte det Ene, Fornuften anerkjender, det
Altomfattende, Altstyrende, Uendelige, hvori den kan finde
Hvile. Og Tankens Modsætning til Folketroens Guder
% emphasis: „hele“ letterspaced (gaps 12, 9, 11 px against 1--5 px solid)
blev til en Modsætning, der omfattede \emph{hele} Menneskelivet.
Saasnart Tanken i sig selv havde fundet sin egen Lov som
Tilværelsens Lov, maatte ogsaa den menneskelige Villie sees
som behersket ved denne Lov: i Fornuftens Fordring, ikke i
en fremmed Villies Bud eller i et Præsteskabs Forskrifter,
maatte Normen søges for den menneskelige Handling. I
Erkjendelse og i Villie søgte da Menneskeaanden det Regulerende
i sig selv, ikke idet den løsnede sig fra Tilværelsens
Heelhed, men idet den begreb sig selv som levende Led i
denne Heelhed, og idet den i sig selv fandt som iboende og
væsensbeslægtet det Guddommelige, der var dens egen og
hele Tilværelsens Grund. Paa denne Grund hvilede den
trygt, uden noget hemmeligt Savn af det, som Tanken bekæmpede,
og hvor den græske Filosofi naaer sit Høidepunkt,
hos Aristoteles, staaer Aanden i Tankens frie Klarhed, i den
trygge Bevidsthed om Eenheden med den guddommelige Fornuft,
der i sin Udfoldning er Verdensaltets Liv, rolig og
fri overfor en Tro, i hvilken den kun seer enkelte Sandhedsglimt
indhyllede i Mythebilledernes brogede Skjul, i
Fantasiens uvilkaarlige eller den beregnende Kløgts vilkaarlige
Tilsløring.

Under vexlende Former, men dog i sit Væsen bestandig
som den samme, træder Modsætningen mellem det frie Menneskelige
og det positivt Religiøse frem i Kristendommens
Tid. Fra Begyndelsen af er Kampen den nye Aabenbarings,
den guddommelige Sandheds Kamp mod den Verdens Viisdom,
i hvilken det Aandelige skal finde en Plads, idet Guds
Viisdom seirer over Menneskets. Her er det Kristelige det,
der strider mod det Sanktionerede, det Bestaaende, det fra
Oldtiden Hellige, idet Verdensviisdommen ogsaa kommer til
at omfatte den Tro, der oprindelig var voxet op af Men%
% --- p. 291 ---
\opage{291}neskelivets Naturgrund, og som nu sees i samme Belysning
som den Tanke, der engang bekæmpede den. Kristendommen
kæmper her, svag i det Ydre, men mægtig i det Indre, fordi det,
den bekæmper, er det Udlevede, hvis Udfoldelsestid er forbi,
hvis Rod i Bevidstheden er udtørret, det, der ikke mere tilfredsstiller
Livets Krav, fordi det kun eensidigt og ufuldstændigt
udtrykker det Menneskelige. Den seirer, idet den
synes at bukke under: af Martyrernes Blod hæver Kirken
sig op som den seirende Kirke, hvis Kors vaier over Romerstatens
Hære, hvis aandelige Magt hæver sig op fra en
jordisk Magts imponerende Grundvold.

I en anden Form træder Kampen frem i Middelalderen.
Mod Kirkens fæstnede Magt, repræsenteret ved et
Præstedømme, udbredt over Landene, samlet i eet Toppunkt
i den Kristi Vikarius, der i Verdens gamle Hovedstad aandelig
hersker over „Byen og Verden“, kæmpes der ikke blot
fra Statens, den verdslige Magts Side, men fra den
Videnskabs Side, som Kirken vil gjøre til sin Tjenerinde
og ligesom indeslutte i en Klostercelle, fjernt fra Naturen og
fjernt fra det naturbestemte Samfundsliv. Tanken, der villigt
og tillidsfuldt har underordnet sig under Kirken, der har
anspændt sin Kraft paa at løse den Opgave, denne havde
stillet den: at knytte Dogmerne sammen til en Kirkelæres
Eenhed, hvis guddommelige Sandhed vel ikke skulde kunne
begribes af den menneskelige Aand, men dog skulde kunne
hævdes af denne som et Ikke-Umuligt; --- denne Tanke
voxede uformærkt under Tjenerforholdet, og idet den voxede,
vaagnede dens Bevidsthed om dens egne Love, dens Trang
til Frihed og dens Trang til Virkelighed. Fra den Gjenstand,
der som det Utilegnelige og Ubegribelige kun paa udvortes
Maade kunde formes, vendte Tanken sig mod det
tilegnelige Virkelighedsindhold, som den kunde øse af en
Natur, der ikke mere skræmmede den tilbage som et Fremmed,
og som den kunde hente op fra Menneskeaandens eget
Dyb. Som Filosofi skilte Videnskaben sig fra Theologien,
den frie Tanke fra den bundne, og Modsætningen, der i den
% --- p. 292 ---
\opage{292}tidligere Middelalder havde kunnet skjules, stod ved Middelalderens
Slutning klar og skarp for Bevidstheden. „En er
Theologiens Sandhed, en Filosofiens“ blev nu Løsenet, og
dermed var Skilsmissen gjort, den verdslige Videnskab
emanciperet. Kirken kunde endnu forkjættre, men dens
Magt over Aanderne var ikke længer den samme, som da
en Keiser barfodet og i Bodsdragt ydmygt maatte trygle en
Pave om Befrielse fra Bannet, og da Tankens Ypperste
maatte tilbagekalde for Concilierne hvad Tanken bekræftede,
Kirken benægtede.

Ved den nyere Tids Begyndelse blev da Kampen en
Kamp mellem den emanciperede Videnskab og den ved
Dogmet bundne, af det Ubegrebne og Ubegribelige beherskede
middelalderlig-scholastiske Theologi, og fra denne svagere
Modstander vendte Videnskaben sig snart mod selve Dogmet.
I den Tid, da Oldtiden gjenfødtes og da Menneskeaanden
begeistret vendte tilbage til de klassiske Forbilleder i Liv og
Digtning, i Kunst og Videnskab, for i dem at søge et Indhold,
den kunde tilegne sig og finde Næring i, blev Kampen
ført med Lidenskab, og i denne Kamp fik Filosofien sine
Martyrer. Snart antog dog Modsætningen en roligere
Skikkelse. Videnskaben behøvede snart ikke længer at kæmpe
for sin Existens; den blev, betrygget ved sin egen Kraft, en
Livsmagt, der tiltvang sig Anerkjendelse, og mod hvilken
% emphasis: „Hele“ letterspaced (gaps 10, 12, 8 px)
den splittede Kirke maatte værge sig. \emph{Hele} den nyere Tids
Filosofi, i dens dobbelte Form som realistisk og idealistisk,
staaer ved sit Princip i Modsætning til det positivt Religiøse,
saaledes som det repræsenteres af Kristendommen, og
% emphasis: „hele“ letterspaced (gaps 12, 10, 9 px)
\emph{hele} den nyere Naturvidenskab er fremmed for denne. Hvad
Kristendommen fornægter: Naturen, er den nyere Videnskabs
Grundlag, og den menneskelige Tanke, der af Kristendommen
% printer's defect: the „e“ of „Synden“ faintly and brokenly printed (raised, partial); transcribed as the word
sees som fordunklet og formørket ved Synden, der har forstyrret
Væsenet, og som kun ledet til Sandhed ved Aabenbaringen,
bliver for den nyere Filosofi en Sandhedskilde,
dens styrende Magt dens egne Love, som den finder stadfæstede
i Naturen. Og ved Videnskabens Modsætning til
% --- p. 293 ---
\opage{293}det positivt Religiøse bestemmes i vor Tid den almindelige
Bevidsthed, direkte eller indirekte. I den moderne Tids
hele Verdens- og Livsanskuelse er Modsætningen udpræget
til den positive Religions Fastholden af en fremmed, overmenneskelig
og overnaturlig Lov og Norm for Erkjendelse
og Handling. For den humane Bevidsthed staaer Naturens
Væsenlighed som Menneskelivets Grundlag fast, og den
tvivler ikke om „den uforvandlede, fuldstændige Menneskenaturs
Evne til, ud fra sit eget Indre at naae til Erkjendelsens
og Handlingens Sandhed, hvis Kilde Menneskevæsenet
gjemmer i sig.“ Og denne humane Bevidsthed, der fjerner
sig fra den positive Religion, finder i sig det idealt Religiøses
levende Kilde ved at gaae tilbage til sit Væsens
Grund, til dets Grundforhold til Tilværelsens evige Magt.

Overalt i de Lande, der gaae i Spidsen for den nyere Tids
Kulturudvikling, finde vi i vor Tid Modsætningen mellem
det Humane og det positivt Religiøse skarpt udpræget og
klart bevidst, og i de mange, for et overfladisk Blik kun
heterogene Former, i hvilke det Humane udpræger sig, er
det ikke vanskeligt at finde de gjennemgaaende Træk, og i de
dybere og alvorligere af disse Former viser sig overalt en
Stræben efter, i selve Menneskenaturens Dyb at finde en
rigere og sandere Erstatning for, hvad det positivt Religiøse
engang ydede Aanden. Stærkere og stærkere hævder det
Humanes Magt sig under Kulturudviklingen, idet Menneskenaturens
Rigdom fyldigere og fyldigere udfolder sig for
Bevidstheden; mere og mere frigjør det sig for Fortidens
hæmmende Baand. Kun i de Lande, der ligge udenfor
Kulturstrømningen eller ved dens Yderkanter, bevarer det
Gamle til en vis Grad sin Betydning for Aanden og sin
Magt over den. Her kan man da see det eiendommelige
Fænomen, at netop denne Fastholden af det Gamle fortolkes
som et Tegn paa, at det er disse Lande, i hvilke den
sande Livskraft er bevaret, at det skal blive dem, der skulle
gaae i Spidsen for Menneskehedens Fremtidshistorie, medens
de store Kulturfolk skulle være udlevede og aandeligt udtømte.
% --- p. 294 ---
% p. 294 opens flush left, without indent: the paragraph of p. 293 continues
\opage{294}Man glemmer at drage en historisk Parallel, der ligger nær.
Da Oldtidens Tro sygnede hen og en ny Livsmagt seirende
hævede sig, da var det ikke paa de Punkter af den gamle
Verden, hvor Tankens Liv var stærkest, Kulturstrømmen
mægtigst, at den gamle Tro gjorde længst og seigest Modstand.
Det var i afsides Egne, i Provindser, der tærede
paa gamle Minder, paa Erindringer om svunden Storhed;
det var i Smaastæder og Landsbyer, at det Gamle bevaredes,
og Navnet Paganisme for det gamle Hedenskab
minder om dette Forhold. Det gjentager sig i vor Tid
med Hensyn til Middelalderens Tro; men nu sætter man
% antiqua: „pagani“ set in roman type, with the roman guillemets »...« (not „...“)
sin Stolthed i at høre til »\textit{pagani}«.

% p. 294: a short centred rule (c. 1.8 cm) between the paragraphs; no section number
\begin{center}\rule{1.8cm}{0.4pt}\end{center}

Under alle de vexlende Former, som Modsætningen
mellem det positivt Religiøse og det Humane har antaget
lige fra den første kristelige Tid og indtil vore Dage, er der
eet Fænomen, der maa tildrage sig Opmærksomheden: det er
den kristelige Livsanskuelses Ubevægelighed under den tilsyneladende
Udvikling. Kun for et flygtigt Blik viser der
sig i den kristelige Livsanskuelse en Udvikling; hvad der antager
Skinnet af en saadan er kun Afspeilingen af Modsætningens
Vexlen, eller det Kristeliges Amalgameren med
Elementer, der ukonsekvent optages fra det Fremmede. I
det Humane viser der sig en Udviklingsmulighed, riig som
dets Væsen, uendelig som dette; i den kristelige Livsanskuelse
er der og maa der være Ubevægelighed, baade paa Grund
af dens Abstraktion, af dens naturløse Aandelighed, og paa
Grund af dens Bundethed ved Autoriteten. Denne Bundethed
er, væsenlig seet, lige stor i Protestantismen og Katholicismen.
Katholicismen har et Fortrin deri, at Traditionen
i den i hvert Øieblik bæres af en levende Individualitet,
paa hvem Tidsaandens Strømninger kunne udøve deres
Indflydelse; men selve Kirkens ufeilbare Overhoved er dog
saaledes bundet ved den Traditionens Sum, der hober sig
% --- p. 295 ---
\opage{295}op bagved ham, at der ikke kan blive Plads for en fri Udvikling.
Det er den samme Traditionens Magt, kun standset
paa et bestemt Punkt: ved den Skrift, der selv kun er
Mindesmærket om en enkelt kristelig Tidsalders Tradition,
der i Protestantismen binder den frie Forskning og den frie
Udvikling. Da Reformationen protesterede mod Katholicismens
Tradition og i Overeensstemmelse med Renaissancetidens
almindelige Stræben, for hvilken den kirkelige Reformation
kun var et enkelt, begrændset Udtryk, tyede tilbage
til det Oprindelige, da var dette Oprindelige ikke det i
Menneskeaanden Primitive: Aandens Frihed i Tanke og
Handling, men det i endelig Tidsbetydning Første: Kristendommens
første historiske Mindesmærker, og man tyede kun
tilbage til dem for til Fordeel for Skrifttraditionen at give
Afkald paa Tankens Ret, paa Handlingens Autonomi, for
at underkaste sig en ydre, blivende Norm, mere bindende end
Katholicismens.

% p. 295: a short centred rule (c. 1.8 cm) between the paragraphs; no section number
\begin{center}\rule{1.8cm}{0.4pt}\end{center}

Naar det hyppigt paastaaes, at den almindelige aandelige
Bevægelse i vor Tid gaaer tilbage mod det positivt
Religiøse, at der gjennemgaaende gjør sig en Trang gjældende
i Aanderne til at gjenoptage dette, saa grunder denne
Paastand sig paa en Forvexling. Man forvexler det Religiøse,
der er et nødvendigt Moment i det fuldstændig opfattede
Menneskelige, et Udtryk for selve Menneskenaturen i
dens Livssammenhæng med det iboende Guddommelige, med
det i de positive Religioners Form Udprægede. Om hiint
vilde hiin Paastand i det Væsenlige være rigtig; om dette
er den usand. Det er hiint religiøse Element, hvilket det
18de Aarhundredes aandløse Materialisme og dens Fornyelse
i vore Dage have villet fortrænge, som nu fra saa
mange Sider, i alle de dybere og ægtere Former for den
humane Aandsretning, stræbes gjenindsat i sin Ret. Idet
Bevidstheden ligesom sænker sig ned i sit eget Dyb og der
mødes med et Høiere, der er Menneskelivets Grund og dets
% --- p. 296 ---
\opage{296}Livsmagt, gjenfinder den i sig det religiøse Forhold, og idet
det er fundet ad Erkjendelsens Vei, kan Menneskeaanden
leve deri, uden at komme i Splid med sig selv og opgive
sig selv og uden at fornægte en Grundvirksomhed af sit
Væsen ved at bringe Viden som et Offer til Troen. Idet
den sædelige Bevidsthed modnes, bliver den religiøs Bevidsthed,
ikke længer „religionsløs Sædelighed“. Men dette Religiøse,
der har sin Spire i det Humane som saadant, er et
væsenlig Andet end den positive Religions Indhold; det er
det almeent, ideelt Religiøse. Derfor blive, samtidig med at
det søges, netop de konfessionelle Differenser, ja selv Differenserne
mellem de forskjellige, historiske Religioner, forsvindende,
% emphasis: „almeent“ (gaps 8--11 px) and „religiøs“ (gaps 8--12 px) letterspaced
og de maae blive det, netop fordi det er det \emph{almeent}
Humane i \emph{religiøs} Form, der søges draget frem og udfoldet.
I vore Dage kan Protestanten og Katholiken, den
Kristne og Jøden, ja den Kristne og Hinduen, der er berørt
af Humanitetens Livsaande, mødes i en Tro, der væsenlig
er een, selv om de ikke have givet denne fælles Tro et
fælles Udtryk. Og denne Tidens Stræben hen imod det
Religiøse i dets Frigjorthed fra de historiske Former --- en
Stræben, der kun ofte mangler fuld Klarhed og Mod til at
vedkjende sig sig selv, --- kan man ogsaa gjenkjende som skjult
under de forskjelligste, tilsyneladende heelt ueensartede Fænomener.
Den er skjult, som forkuet og uklar, under den
halvspekulative, halvkristelige Theologis Forsøg paa en halvveis
Rationaliseren af Dogmet, paa en halvveis Humaniseren
af det Kristelige, ved Hjælp af en halvveis Opgiven af
det principielt Kristelige. --- Den er skjult i den almindelige
Theologis Forsøg paa at gaae bagved de historiske, konfessionelle
Udtryk for Kristendommen for at finde en oprindelig,
ideal Kristendom, der skal have foresvævet Stifterens Tanke,
men aldrig have fundet et tilsvarende Virkelighedsudtryk; thi
denne ideale, uhistoriske Kristendom bliver i Virkeligheden
kun et Udtryk for det abstrakt almene Element i de positive
Religioner, i hvilket de falde sammen med Fornufttroen. ---
Den er som skjult virkende tilstede i den almindelige Opfat%
% --- p. 297 ---
\opage{297}telse af de positive Religionsformer, der ere blevne til i
vore Dage. Selv de, der ikke føle Hvile i de ældre positive
Religioner, betragte dog de nye med Ligegyldighed eller
Mistillid. Vor Tids Bevidstheds Forhold til den positive
Religion bliver det samme som dens Forhold til Miraklet.
Ligesom dette bliver den positive Religion betragtet som noget
Forbigangent; selv om man ikke nægter dens Realitet overhovedet,
opfatter man den som Noget, hvis Tilblivelse kun
hører Fortiden til, kun der har Virkelighed, ikke i det Nærværende.
Men denne Opfattelse er kun den indirekte Benægtelse
af dens Realitet overhovedet. --- Og endelig kan hiin
Stræben ogsaa være skjult under selve Tilslutningen til de
absurde, positive Religionsformer, som vor Tid har skabt;
thi hyppigt er det netop en uklar Søgen efter det frie
Religiøse, der bringer til at kaste de tidligere positive
Former bort, om end kun for at finde en midlertidig Hvile
i en ny positiv Religion.

% p. 297: a short centred rule (c. 1.8 cm) between the paragraphs; no section number
\begin{center}\rule{1.8cm}{0.4pt}\end{center}

Det almeent Religiøse, til hvilket den humane Bevidsthed
i vor Tid søger tilbage, hvor den opfatter sig selv
dybere og fuldstændigere, denne „Menneskehedens Religion“,
der er Udtrykket for Menneskevæsenets oprindelige ideelle
religiøse Trang, bliver klaret for Bevidstheden gjennem Viden,
gjennem den sande, gjennemførte Erkjendelse, gjennem den
filosofiske Tænkning. Men dets Forhold til denne Erkjendelse
sætter ingen blivende Sondring mellem Tænkerens
Religion og Mængdens Religion. Hvad der er begrundet i
det almeent Menneskelige, maa kunne udformes i almeent
tilgængelige Bevidsthedsformer. Ikke i enhver Form vil den
høieste Sandhed, som Videnskaben udklarer, være tilgængelig
for Alle, men i en eller anden Form vil den kunne være
det. „Sandhedsmonopoler til Fordeel for en enkelt Menneskeklasse
vilde være Forurettelser mod alle de andre og virkelige
Injurier mod Menneskeheden“ (Lichtenberg). Sandheden
% --- p. 298 ---
\opage{298}er det Forenende, det, der skal blive Alles Eiendom, et Almeenlivs
Medium. Og det Religiøse, der har sin Rod i
selve Menneskevæsenets Grund, det frie og frigjorte Religiøse,
vil, ligesom det kan blive Alles Eiendom, tilfredsstille
% emphasis: „hele“ letterspaced (gaps 11, 11, 10 px)
\emph{hele} den menneskelige Aand, tilfredsstille ogsaa Hjertets
Trang. Det vil kunne det, fordi det er selve Menneskevæsenets
Udtryk, og fordi det for den religiøse Bevidsthed
Væsenlige har Evnen til at udpræge sig i de forskjelligste
Former efter den fremskridende menneskelige Udviklings forskjellige
Stadier.

% p. 298: a short centred rule (c. 1.8 cm) between the paragraphs; no section number
\begin{center}\rule{1.8cm}{0.4pt}\end{center}

For den, der med et uhildet og omfattende Blik betragter
den religiøse Bevægelse i vor Tids Bevidsthed, vil
der ikke let kunne opstaae en Tvivl om dens virkelige Retning
og om dens Fremtidsmaal. Men der ligger i de
historiske Forhold, under hvilke den nyere Tid har udviklet
sig, Meget, der kan gjøre det forklarligt, at Opfattelsen af
Fænomenerne saa hyppigt forvirres, hvor Usikkerhed med
Hensyn til det Principielle virker sammen med Tendenser,
der fordunkle Erkjendelsen. Hvis det Kristelige havde været
% emphasis: „ny“ letterspaced (n--y gap 11 px against 2--3 px in the solid „en“ before it)
det historisk Oprindelige, hvis det var kommet ind i en \emph{ny}
% p. 298: a faint grey speck after „Aanderne,“ (show-through or an inked space), not a character
Verden, og hvis det ene havde kunnet udfylde Aanderne, saa
havde Opfattelsen af det kunnet være sikker og let: dets
virkelige Natur vilde ligesom have paatvunget sig Erkjendelsen.
Men da Kristendommen kom ind i Verden, da vare
allerede store Kulturperioder afsluttede, eiendommelige Grundformer
af det Humane vare allerede udfoldede og udtømte.
Og Kristendommen har, som det for Verden Fremmede, som det
Naturen Fornægtende, altid kun kunnet udfylde eet Gebeet
af Aanden, en enkelt Trang, en enkelt Livsform, ikke omfatte
Menneskevæsenet efter dets hele Omfang. Den har
bestandig havt ved Siden af sig et andet Element, et andet
Livsprincip, og idet dette udviklede sig samtidig med Kristendommens
Bestaaen og under Vexelvirkning med det
Kristelige, kunde den Forvexling let opstaae, at det, der kun
% --- p. 299 ---
\opage{299}var blevet til paa samme Tid som det Kristelige, tildeels i
Kamp imod det, kunde sees som fremgaaet af selve det Kristeliges
Princip, som dettes rige Udfoldelse. Derved blev
det muligt, at det, der kaldes Kristendom, kunde blive et af
saa ueensartede Elementer sammensat Fænomen, at de, der
staae indenfor Kristendommen, men have modtaget deres
religiøse Bevidstheds Grundindhold, dens Støttepunkter, i
denne Sammenvoxethed med det Fremmede, kunde blive
usikkre over det Kristeliges Grændser mod det Humane, og
at selv de videnskabeligt Dannede blandt Theologiens Repræsentanter
kunde betragte en skarp principiel Udsondring af
det specifisk Kristelige, hvad enten den fandt Sted fra et
Standpunkt indenfor eller udenfor Kristendommen, som en
„uforsvarlig Misforstaaelse“, medens de selv fremsatte en
Opfattelse af det Kristelige, der, idet den med en skrøbelig
Dialektik søgte at annektere det Humane, gjorde det Kristelige
til et Udvisket og Uklart, og kun skjulte Principløsheden
ved Frasernes Taageskyer. ---

Vi skulle i det Efterfølgende forsøge at lade Repræsentanter
for de to i vore Dage stridende Opfattelser af Kristendommen
komme til Orde: for den, der sætter Modsætningen
mellem det Kristelige og det Humane, og for den,
der lader det første omslutte det sidste. Og idet vi saa ad
Historiens og Begrebsudviklingens Vei begrunde vor Tilslutning
til den Opfattelse, der fastholder Modsætningen,
ville vi sammenstille det kristeligt Religiøse med Humanitetens
Religiøsitet, der for os staaer som den religiøse Bevidstheds
Maal. For denne ville vi da søge et Udtryk, ikke
udelukkende hos nogen enkelt af dens Talsmænd, men i
det, der i de forskjellige Nuancer, i hvilke den har givet sig
Skikkelse, viser sig som et Fælles og Gjennemgaaende, som
det, hvori dens væsenlige Motiver blive klare.

% p. 299: signature set right in heavy (bold) Fraktur; then a centred double rule (c. 2.6 cm) closes installment I; below it, faint show-through only
\begin{flushright}\textbf{Constans.}\end{flushright}
\begin{center}\rule{2.6cm}{0.4pt}\\[-10pt]\rule{2.6cm}{0.4pt}\end{center}

\clearpage
% ===== II. (pp. 433--456) =====
% --- p. 433 ---
% p. 433: the upper part of the page is blank (the installment opens a new page; no preceding article, no running head; folio 433 at the head; show-through only); the head is set in a display blackletter (textura) in two lines, „Det Kristelige og det Humane som Modsætninger“ / „i vor Tids Bevidsthed.“; a short rule below it; then the numeral „II.“ in roman type, centred
\opage{433}\begin{center}\textbf{\large Det Kristelige og det Humane som Modsætninger i vor Tids Bevidsthed.}\end{center}

\secnum{II.}

% p. 433: the opening „I“ is a slightly enlarged Fraktur initial, standing a little below the line
I Opfattelsen af det Kristelige staaer der i vore Dage,
indenfor selve Kristendommens Omraade, to Modsætninger
ligeoverfor hinanden, skarpt sondrede i Principet. Den ene
har fundet sit klassiske Udtryk i vor Literatur gjennem en
% emphasis: „Søren Kierke-|gaard“ letterspaced over the line break (gaps in „Søren“ 15, 9, 10, 9 px; „Kierke“ 10--15 px; „gaard“ 6--11 px, against 1--5 px solid)
eminent begavet Personlighed: gjennem \emph{Søren Kierkegaard}.
For ham staaer det Kristelige som det, der er
ubetinget ueensartet med alt det Verdslige og med det for
Tanken Tilgængelige. Kristendommens Centrum: Gudmennesket,
er det absolut Ubegribelige, det ubetingede
% p. 433: a faint grey speck after „i“ (lighter than the type; show-through or an inked space), not a character
Paradox, til hvis Anerkjendelse i Troen ikke Fornuftens
Slutninger, men Syndsbevidsthedens Fortvivlelse bringer.
% p. 433: „Det er det;“ -- a semicolon, as printed
Det er det; hvortil der naaes ikke ved Tanken, men tiltrods
for Tanken og gjennem et Brud med den. Troen
griber det Ubegribelige, og idet den med Inderlighedens
Lidenskab fastholder det, værner den netop med al sin Kraft
og al sin Lidenskab om dets Ubegribelighed mod ethvert
Forsøg af Tanken paa at tilegne sig det. Og ligesom den
kristelige Tro af Kierkegaard kun opfattes som fremgaaet af
et Brud med Tanken, saaledes opfattes ogsaa det kristelige
Liv, der spirer frem af denne Tro, som fremgaaet af et
% --- p. 434 ---
\opage{434}Brud med Verden. Kristendommen er Lidelsens Religion,
dens Bud er det at afdøe fra Verden; kun gjennem denne
Død viser den Veien til Livet. Og dette negative Forhold
til Verden er ikke indskrænket til en enkelt Tid, som f. Ex.
til Kristendommens første historiske Fremtræden. Saalidt
som Troens ubetingede Paradox nogensinde bliver begribeligere
ved Tiden, saalidt bliver det Liv, der hviler i det, nogensinde
eensartet med Verdens Liv. Kristendommen, der ved
at fordre Døden, Jeghedens Død og Forsagelsen af alt det
Verdslige udenfor og i Individet, vækker Selviskhedens og
Verdslighedens Modstand, maa bestandig staae lige fjendtlig
overfor „Verden“, den maa til enhver Tid blive bespottet,
forhaanet, hadet, forfulgt. Og hvor Historien viser et fredeligt
Forhold, der er det netop Tegnet paa, at Kristendommen
har tabt sin eiendommelige Karakteer, at den er forvansket
ved en Blanding med det Ueensartede, at den har nedstemt
den ideale Fordring, fornægtet sig selv. I vor Tids Kristenhed
gjenkjender derfor Kierkegaard ikke Kristendommen; han
seer i den kun et gjennem mange Aarhundreder fortsat, omhyggeligt
tilsløret Bedrag; i den af Staten beskyttede Kirke
seer han kun et Kristendommens Vrængebillede, hvis Forkyndere
leve af Kristus istedetfor at leve for ham. Han
opkaster tvivlende det Spørgsmaal: om Kristendommen endnu
findes paa Jorden, og idet han med sit Ideal af det Kristelige
for Øie gaaer tilbage gjennem Kirkens Historie ligetil dens
første Tid, bringer han hiin Tvivl med sig ligetil denne
og spørger tungsindig: har der vel nogensinde været Kristne?
er Kristendommens ideale Fordring nogensinde tilfulde blevet
realiseret? Og dette Spørgsmaal er konsekvent; thi kan vel
i Verden den Religion finde sit fuldkomne Udtryk, hvis
% emphasis: „ikke“ letterspaced
Princip er, \emph{ikke} at leve i Verden, at være død for Verden?

I denne Opfattelse af det Kristelige er paa een Gang
den skarpeste Modsætning udtrykt til „Kristenhedens“ Kristendom
og til det Humane. I Modsætning til Begrebet om en
Kristenhed, i hvilken Individet ligesom er Kristen ved Fødselen
og vedbliver at være det i Kraft af Sammenhængen
% --- p. 435 ---
\opage{435}med et objektivt Bestaaende, akcentuerer Kierkegaard Begrebet
om „den Enkelte“. Ansigt til Ansigt vil han stille Individet
som „Enkelt“ med Gud, med det Religiøses uendelige Fordring,
for at lade Afgjørelsen skee i Inderlighedens Dyb, i
det Møde, hvor det Guddommelige og Menneskelige berøres
i Selvets Inderste. Enhver objektiv Vei til Gudsforholdet
afskjærer han, ikke blot den, der gaaer gjennem den objektive
Tænkning, men ogsaa den, der gaaer gjennem Kirken. Det
subjektive, personlige Forhold til Sandheden bliver, ifølge
Menneskets Existensbetingelser, det ene sande: i uendelig
Bekymring for sig selv, i den Bekymring, som Syndsbevidstheden
vækker, skal Individet i Troens Subjektivitetsforhold,
der er en uendelig lidenskabelig Interesses Forhold, som Enkelt
vinde Forholdet til det Guddommelige saaledes, som det
er aabenbaret som „Guden i Tiden“, som det Paradox, der
for Tanken er det Absurde, for Troen det uendeligt Visse,
der i hvert Øieblik i Frygt og Bæven vindes ud af den
objektive Uvished.

„Den Enkelte“ er Udgangspunktet og maatte være det
for den, der i „Kristenheden“ saae en Forvanskning, som paa
een Gang træder frem i Sammenblandingen af det Kristelige
% UNSURE: printer's defect: the second „e“ of „Individets“ prints without the closing stroke of its eye, like a Fraktur „c“ (a broken e or a wrong sort); transcribed as the word
og det Verdslige og i Individets Absorption i det kirkelige
Samfund, det upersonlige Almene. Men „den Enkelte“ er
ikke, som det endnu ganske nylig ved en iøinefaldende Misforstaaelse
er blevet paastaaet, Slutningspunktet for Kierkegaard.
Først maatte Blændværket brydes, og det maatte
gjøres klart for Individet, at dets Vei til Gud kun kan
gaae gjennem Inderlighedens Snevring, ikke ad det Objektives
Alfarvei; kun paa denne Betingelse kunde Omvæltningen
i Sjælene, som Kierkegaard satte sit Liv ind paa, blive til
Virkelighed. Men naar de „Enkelte“ hver for sig vare
traadte ind i det sande Gudsforhold, saa fik Menighedens,
Samfundets Begreb atter sin Gyldighed, fordi de Enkelte
mødtes i det samme evige Maal. Der blev et Samfundsforhold
for de Enkelte, men saaledes, at Gud var Mellemstadiet
mellem den Enkelte og Samfundet, ikke Samfundet
% --- p. 436 ---
\opage{436}Mellemstadiet mellem den Enkelte og Gud. Dette Forhold
er saa klart udtalt i Kierkegaards Skrifter, at kun en Trang
til at bære en Overlegenhed tilskue ved at faae Kierkegaards
% p. 436: „ethisk-kristeligt-|kirkeligt-humant“ -- the line breaks after a compound hyphen, which is kept
Standpunkt indordnet som Moment i et høiere ethisk-kristeligt-kirkeligt-humant
kan forklare den Misforstaaelse, der tillægger
Kierkegaard en anden Opfattelse.

Men det Liv, der sættes som Opgave for den Enkelte,
medens han staaer ene Ansigt til Ansigt med Gud, og medens
han i sit Gudsforhold finder Forholdet til de andre Individer,
bliver et Liv, der skarpt skiller sig fra det Humane.
Hvad der bærer det, er det, der er utilgængeligt for Tanken:
det ubetingede Paradox, det, der for Tanken er det Absurde.
Men Paradoxet, der strider mod alle Tankens Love og
Former, idet det udsiger Foreningen af Gud og det enkelte
Menneske i Kristi Person, det Eviges Tilbliven i et Moment
af Tiden, er for den Troende den høieste Sandhed. Det
Sande dømmer da det, der er i Strid mod det, som det
Usande, og der bliver ikke i Aanden en Plads for et Tankens
Liv; der bliver ingen, end ikke en relativ, Gyldighed
for Tanken. Paradoxet er med Rette blevet betegnet som
„Underet i Tankens Verden“, og det er med Rette blevet
fremhævet, at ligesom Underet i Naturens Verden ophæver
Naturens og Naturlovens Begreb, saaledes maa Underet i
Tankens Verden omstyrte Tankens Love og lade al Videns
Vished opløse sig i Tvivlen. Fordringen til det troende
Individ maatte konsekvent blive, „at det i sit Tankeliv skal
kaste al Verdens Virkelighedsindhold bort for blot at være med
Syndsbevidsthedens Paradox i det paradoxe Gudsforhold.“
Dette maa blive Fordringen, fordi Paradoxet ikke blot modsiger
% emphasis: „al“ letterspaced (a--l gap 11 px against 4--6 px in the solid words about it)
en enkelt Form af den menneskelige Tænkning, men \emph{al} menneskelig
Tænkning; fordi der her ikke er en Strid mellem to
Former af det Rationelle, men mellem det Rationelle i sin
Heelhed og det Antirationelle. Det Absurde, Kierkegaards
% emphasis: „al“ letterspaced (a--l gap 10 px)
% footnote *) on p. 436
kategoriske Bestemmelse for Paradoxet, udelukker \emph{al} Tanke.\footnote{Jfr. Nielsen: Religionsfilosofi, S. 8, 34, 301 ff. o. fl. St.}

% --- p. 437 ---
\opage{437}Men ligesom den Troendes Liv, saaledes som Kierkegaard
opfatter det, danner Modsætningen til det humane Liv
ved at fornægte Tanken, der udtrykker en Væsensbestemmelse
i det Menneskelige, saaledes kommer det i Strid med det,
fordi det konsekvent ikke kan give Plads for det konkret
Sædelige, for alle de Former af et ethisk Liv, der hvile paa
en Naturgrund i det Menneskelige.

Efter Kierkegaards Opfattelse bliver det religiøse Forhold
ikke blot det høieste Forhold, til hvilket det ethiske Forhold
viser hen, idet det viser ud over sig selv, men det bliver i
den Forstand det absolute Forhold, at det ikke kan „udtømme
sig eller blive konkret“ i de ethiske relative Forhold. Det
bliver da i den Forstand det absolute Forhold, at det for sig
gjør Krav paa al Menneskets Kraft og gjør Krav paa den
i hvert Øieblik, thi kun idet jeg altid, i hvert Øieblik „forholder
mig absolut til det Absolute“, bevarer jeg Forholdet
efter dets Væsensbestemthed. Men der bliver da ingen Kraft
tilovers til de ethiske Forhold: de maae forsvinde for det Absolute,
Uendelige, netop fordi de falde udenfor dette; ligesom
Skabningens endelige Kraft maa forsvinde, naar den stilles
ligeoverfor Almagtens uendelige Magt. Der gjøres derfor
ogsaa den Fordring til Mennesket, at det til Fordeel for det
% emphasis: „alle“ letterspaced
uendelige religiøse Forhold skal kunne forsage \emph{alle} de ethiske
Forhold. Det Religiøse bliver ikke Sjælen i det Ethiske,
gjennemtrænger ikke de ethiske Forhold saaledes, at det, idet
det hæver dem op i en høiere Uendelighedens Sfære, selv faaer
Livsvirkelighed derigjennem, men det staaer overfor dem som
det Uendelige, der fortærer det Endelige, som det Guddommelige,
der bringer Døden til det Menneskelige, hvis Øie har seet det.
I Kristendommens Fordring til at afdøe, til at hade sit naturlige
Liv, ligger det indesluttet, at den naturlige Tilværelses
Indhold forsvinder som et Intet i sin ubetingede Modsætning
til den evige Sandhed, og med dette Naturgrundlag forsvinder
Betingelsen for den Udfoldelse af det Ethiske, ved
hvilken det alene kan faae Virkelighed. Jo bestemtere Kierkegaard
udpræger sit Standpunkt under Oppositionen mod det
% --- p. 438 ---
\opage{438}Bestaaendes Principløshed, desto bestemtere træder ogsaa denne
Konsekvens frem. I Polemiken imod Ægteskabet i „Øieblikket“
er Konsekvensen ufordulgt udtalt.

% p. 438: a short centred rule (c. 1.8 cm) between the paragraphs; no section number
\begin{center}\rule{1.8cm}{0.4pt}\end{center}

Kontrasten til den Kierkegaardske Opfattelse af det Kristelige
danner nu den, der udstrækker det Kristeliges Navn til
alle de Former af aandeligt Liv, som have udviklet sig samtidig
med og under Paavirkning af Kristendommen. Den
seer ikke Kristendommen som Modsætningen til vor Tids
aandelige Liv, men som dets besjælende Princip; den taler
om en kristelig Stat, en kristelig Ethik, en kristelig Kunst,
en kristelig Videnskab, med eet Ord: den anvender Kristendommens
Navn paa alle Former af det Humane; den kalder
den Verden, i hvilken Kristendommen er som en fremmed
Gjæst, en kristelig Verden. Den skarpe Sondring mellem
Kristendommens og Verdens Princip lader den, om den end
bestandig viser hen til det Apostoliske som det Orienterende,
kun være berettiget i Kristendommens første Tid, i hvilken
Kampen med den gamle Verdens Former maatte føre til et
exklusivt Forhold. Fra det Øieblik, da Kristendommen blev
seirende, skulde derimod dens Opgave ikke være at sondre
sig fra Verden, at afdøe fra den, men at gjennemtrænge
den; dens guddommelige Magt skal netop vise sig deri, at
den har formaaet dette, og dens Mulighed til som Livsprincip
at gjennemtrænge det Ueensartede søges deri, at Kristendommen
er en ny Skabelse, der virker forhøiende, „potenserende“,
tilbage paa Naturen og paa de aandelige Kræfter og
Love. Ikke Flugt fra Verden, men Fordybelse i den, ikke
Kamp mod den, men Fred med den og Herredømme over
den bliver da Kristendommens Maal, der i den kristelige Stat
naaes saaledes, som det under Endelighedens Vilkaar kan
naaes, og Statens Anerkjendelse af det Kristelige, dens Begunstigelse
% p. 438: a faint speck under the „i“ (show-through), not a character
af det i Retning af verdslige Goder, bliver et nyt
Beviis for, at den med Verden forsonede, humane Kristendom
er den sande. Og den samme Anerkjendelse yder ogsaa
% --- p. 439 ---
\opage{439}Fornuften Kristendommen. Om end den ugjenfødte, formørkede
Fornuft ikke begriber Kristendommens Mysterier, saa er
dog den troende, ved det nye Livsprincip potenserede Fornuft
istand til at trænge ind i deres Dybder, og om den end ikke
naaer til logisk eller mathematisk Klarhed, kan den dog udforske
Troessætningernes Sammenhæng og derved deres
Grund og udvikle en Troesvidenskab, der er al Videnskabs
Toppunkt og Norm. Den kristelige Filosofi hæver sig, idet
Tanken laaner Vinger fra Fantasien, til i Kristendommens
Gud at øine en evig Natur, der er renset fra al den grovere,
ved Menneskets Synd besmittede Naturs Ufuldkommenhed.
Den seer i denne evige Natur i Gud Muligheden til at
hævde det Naturliges relative Berettigelse i det Menneskelige,
til en Forsoning af Modsætningerne i Menneskelivet. Og
i det Guddommeliges Forening med det Menneskelige i
Kristus seer den det høieste Udtryk for den Forsoning, Tanken
fordrer: mellem den uendelige Aand og den endelige, mellem
Skaberen og den skabte, aandelig-naturlige Tilværelse. Saaledes
sees Kristendommen som forsonet med Verden og med
Tanken, som deres nye, fuldendende Livsprincip.

For denne Opfattelse bliver den Kierkegaardske en fordømmelig
Eensidighed, ligesaafuldt som den „emanciperede“
Fornufts Hævden af det Humane som det Sande i dets
Sondring fra det Kristelige. I Sondringens Sted træder
overalt Mæglingen; den rige „høiere Eenhed“ gjør sig med
værdig Overlegenhed gjældende overfor de fattige Eensidigheder;
det Humane optages nedladende som tjenende Moment,
og med rhetorisk Pathos forkyndes Kristendommens Triumfer.

For et flygtigt Blik kan Alt dette være saare skjønt at
skue, og den Trang til en Forsoning, vunden uden Villiens
Kamp og Tankens Anstrengelse, der er saa vidt udbredt, gjør
det psykologisk let forklarligt, at en saadan Udjevningstheori
finder Bifald hos den store Mængde, der ikke trænger til
Klarhed, men til Beroligelse, og som lettere modtager den
glatte Frase end den skarpe Tanke. Seer man lidt nøiere
paa Sagen, saa opdager man dog uden Vanskelighed, at den
% --- p. 440 ---
\opage{440}glatte Frase kun sammenbinder løst sammenheftede Modsætninger,
at den høiere, rigere Eenhed, der skal være vunden,
kun er tilveiebragt ved en saare fattig, overalt stereotypt
gjentaget Methode, der til Resultat kun kan give Skin, ikke
Virkelighed, Fraser, ikke Begreber. Den stereotype Methode
er simpelthen den, at man tager de Modsætninger, man vil
bringe sammen, og forbinder dem i een Sætning som Subjekt
og Prædikat, og at man saa sender den saaledes dannede Frase
ud i Verden med den Vedtegning, at den „maa tænkes“, hvad
der vel nærmest maa forstaaes ligesom det: „maa trykkes“
% antiqua: „imprimatur“ set in roman type (within parentheses)
(\textit{imprimatur}), som i sin Tid Censor satte paa de Bøger, der
sendtes ud i Verden. Til Mere end en Tilladelse til at tænke
Frasen kan man ialfald ikke godt naae; en virkelig Tænken af
den bliver en ligefrem Umulighed. Idet Modsætningerne begge
finde deres Udtryk i Frasen, maae disse deres Udtryk ligesom
neutralisere hinanden: Prædikatet, der føies til det ueensartede
Subjekt, maa, istedetfor at bestemme Subjektet, berøve det
dets Indhold og kun lade en ubestemt, vag Forestilling tilbage
for Bevidstheden, en Forestilling om et Intet, og naar
saa Modsætningerne i deres Bestemthed atter trænge sig frem
for Bevidstheden, er Modsigelsen og Spliden der paany, og
Eenheden maa forsvinde, naar ikke Bevidstheden kan dysses i
en blid Slummer ved Frasens monotone Musik. Et enkelt
Exempel vil være nok til at anskueliggjøre det Sagte. Modsætningen
mellem den positive Religion og Filosofien i Opfattelsen
af det Guddommeliges Forhold til Tilværelsens Væren finder sit
prægnante Udtryk i Modsætningen mellem Skabelsesforestillingen
og Begrebet om en Verdensudfoldelse, en Kosmogoni. For
Religionen er Tilværelsen et Værk af en skabende Almagts ubegribelige
Vilkaarlighed; for Filosofien har Tilværelsen Væren
ifølge selve det Guddommeliges evige Væsensnødvendighed; den
er en Udfoldelse af selve det Guddommeliges Væsensbestemmelser.
Her staae skarpe, som det skulde synes uforenelige,
Modsætninger overfor hinanden; men hiin Methode forener
dem let --- i en Frase, der „maa tænkes“. Først bliver nemlig
Skabelsesforestillingen rationaliseret ved Kosmogoniens Begreb:
% --- p. 441 ---
\opage{441}det skal ikke være en som naturløs Aand bestemt Gud, som
ved en ubegribelig Almagtens Handling af Intet frembringer
en Tilværelse, der som materiel ikke har noget Tilknytningspunkt
i det Guddommeliges Væsen og derfor ikke kan begribes
som Virkning af det. Det skal være en Gud, der har en
Natur i sig, som skaber af fri Nødvendighed, og det skal
ligge i Skabelsens Begreb, at Gud ikke frembringer et Dødt,
men et Levende, at han frembringer en Skabning, der tillige
frembringer og udvikler sig selv i en given Selvstændighed.
% emphasis: „maa“ (gaps 12, 9 px) and „tænkes“ (gaps 8--10 px) letterspaced
Derfor \emph{maa} Skabelsen \emph{tænkes} at begrunde Kosmogonien
eller Verdens Selvudvikling, dens Genesis. Men naar saa
dette Skridt henimod Naturalismen er gjort, naar det, der
for den religiøse Bevidsthed gjør Skabelsen til Skabelse, er
udvisket, saa bliver Foden forsigtigt trukket tilbage, og der
fortsættes: „Verdens kosmogoniske eller naturlige Begyndelse
er den relative, den endelige, som derfor ogsaa er splittet i
sporadisk Mangfoldighed. Al Organisation danner sig sporadisk,
og betragtet fra det kosmogoniske eller det naturlige
Synspunkt kan Verden siges at have utallig mange Begyndelser,
at begynde fra uendelig mange Livsspirer, der, forsaavidt
som det naturlige Synspunkt udelukkende fastholdes,
kun have deres fælles Eenhed i Kaos. Men de uendelig
mange naturlige Begyndelser ere alle begrundede i den ene
skaberiske, overnaturlige Begyndelse, i den guddommelige Logosvillie,
der i sig indeslutter Livets og Lysets Væld og af sit
aandelige Dyb løslader den hele Mangfoldighed af Livskræfter
% p. 441: a faint speck between „Livskræfter“ and „til“ (show-through), not a character
til Selvbevægelse.“ Altsaa: Skabelsen var kosmogonisk,
det Overnaturlige naturligt, og nu bliver til Gjengjæld
Kosmogonien Skabelse, det Naturlige overnaturligt. Saaledes
maa Eenheden tænkes; ja, det vilde være en uforsvarlig
Eensidighed at adskille Bestemmelserne. Men ligesom den
Skabelse, der var kosmogonisk, ved dette Prædikat netop mistede
det, der skulde gjøre den til Skabelse, saaledes mister nu den
% emphasis: „begrundet“ letterspaced (gaps 10--13 px)
Kosmogoni, der som \emph{begrundet} i Skabelsen selv bliver
Skabelse --- idet den relative, naturlige Begyndelse forsvinder
i den overnaturlige, det Materielle i det „aandelige Dyb“ --- til
% --- p. 442 ---
\opage{442}Gjengjæld netop det, der gjør den til Kosmogoni. Vi faae
da tilbage to tomme Forestillinger, to Navne uden alt Indhold,
som vi have Frihed til at tænke, hvis vi ellers kunne
tænke Noget ved det Indholdsløse; og hvis vi ikke kunne det,
og vi for ikke at lade vor Tanke tabe sig i det Tomme kalde
Begrebernes Indhold tilbage og saa forsøge at benytte Tilladelsen
til at tænke deres Eenhed, saa splittes Modsætningerne
atter ad. Saa staaer Skabelse atter imod Kosmogoni, Vilkaarlighed
imod Nødvendighed, Fantasibilledet af Naturen
i Guds „aandelige Dyb“ imod den reale, virkelige, materielle
Natur, Underet imod den væsensgyldige Naturs lovmæssige
Udvikling, det Overnaturlige mod det Naturlige. Og Modsigelsen
hæves ikke ved saadanne Logogryfer som, at Gud
„indskrænker“ sin Almagt, men „ikke bortgiver den“, at han
„frembringer en Skabning“, der tillige „frembringer sig selv“,
en Skabning, der besidder en „given Selvstændighed“. Eenheden
fordunster i de tomme Talemaader: begge Modsætningerne
forvanskes, og der kan kun bringes Mening i Fraserne,
naar man istedetfor det halvt Rationelle, det halvt
Naturmæssige sætter det heelt Antirationelle, det heelt Supranaturalistiske,
tager det, der hører til den ene Modsætning,
som en blot Talemaade, der ikke skal betyde videre, og lægger
alt Eftertrykket paa det, der hører til den anden. Men saa
er Mæglingen opgivet, og vi ere vendte tilbage til den Opfattelse,
der fordømtes som en Eensidighed.

% p. 442: a short centred rule (c. 1.8 cm) between the paragraphs; no section number
\begin{center}\rule{1.8cm}{0.4pt}\end{center}

Vi kunde indskrænke os til at paavise det Skrøbelige i
Mæglingstheoriens Methode; men for at give et fuldstændigt
% antiqua: „Juste-milieu-“ set in roman type, „Theorien“ in Fraktur
Materiale til Dommen over \textit{Juste-milieu}-Theorien gaae
vi tillige Historiens og Begrebsudviklingens Vei for ad den
at sikkre os Forstaaelsen af det Kristelige, Opfattelsen af
dets væsenlige og principielle Forhold til det Humane.

Naar vi ville finde de historiske Forudsætninger for
Kristendommens Tilblivelse, blive de væsenligt at søge i den
% p. 442: a faint mark before „Aands“ (show-through), not a character
græske og i den jødiske Aands Udvikling. Fra forskjellige
% --- p. 443 ---
\opage{443}Udgangspunkter stræbe disse Folkeaander, Repræsentanterne
for den ariske og den semitiske Stamme, hen mod eet Udviklingsmaal,
og idet de mødes, ere Betingelserne givne for
Kristendommens Tilbliven. Hvad der i den kan være optaget
af andre, østerlandske, Elementer, er i ethvert Tilfælde
kommet til den gjennem det Græske eller det Jødiske.

% emphasis: „græske“ letterspaced (gaps 6--10 px against 4--5 px in the solid „Aands“ beside it)
Betragte vi da først den \emph{græske} Aands Udvikling, saa
bliver her af væsenlig Betydning Dualismens Frembryden
i den græske Tænkning og Livsanskuelse. Det græske Folk,
Skjønhedslivets Repræsentant i Verdenshistorien, havde fundet
sit høieste Udtryk for Anskuelsen af Tilværelsen i Opfattelsen
af den som en Kosmos, som et Verdenskunstværk; det havde
fundet sit høieste Udtryk for det ideale Menneskeliv i Opfattelsen
af det som et sædeligt Kunstværk, som en harmonisk
% p. 443: a faint vertical stroke before „Fornuften“ (an inked space or rule edge), not a character
Formen af Natursiden i Mennesket ved Fornuften. Skjønhedsanskuelsen
var saaledes for Grækerne det Høieste; den
gjød som en Straaleglands over Tilværelsen og Menneskelivet,
der saaes som hvilende i Harmoniens Ro. Og i Forholdet
til Guderne traadte den samme Skjønhedsanskuelse
i Forgrunden. Naar en Græker stod overfor Fidias’s
Billede af den olympiske Zeus eller af Pallas Athene, da
var hans Andagt den, at lade det ideale Billedes Skjønhedsharmoni
% printer's error: „Stihed“ (for „Stilhed“; no l, confirmed at the image); transcribed as printed
sænke sig over Sjælen, bringe Stihed, Ro og
Harmoni ind i dens Attraa og Lidenskab, at lade al Smerte
dulmes ved Anskuelsen af den guddommelige, nærværende
Skjønhed.

I denne Skjønhedsanskuelse af Tilværelsen er det Ideales
Herredømme hævdet, Naturen er underlagt Aandens Magt.
Men Skjønhedsharmoniens Eenhed af Aand og Natur er dog
kun en svigefuld Eenhed, svigefuld, fordi Kunstværket i sin
umiddelbare Eenhed for Anskuelsen forbinder to Modsætninger:
Stoffet og Formen, der staae som oprindeligt forskjellige,
som ikke oprundne af hinanden. Naar den græske Tanke
naaer sin fulde Udvikling, naar den ligesom vikler Ideen ud
af Sandseverdenens skjulende Slør og vil begribe dens Væsen
og dens Forhold til det Stof, der ved den bliver til et
% --- p. 444 ---
\opage{444}Verdenskunstværk, saa skilles Modsætningerne ad for Tanken,
og den græske Tanke mægter ikke fra Sondringen at føre
dem tilbage til en Eenhed, fastere og sikkrere end Harmoniens
Skjønhedseenhed, fordi den som græsk Tanke har sin Grændse
i Skjønhedsanskuelsen, er bunden ved denne som ved en
uoverstigelig Skranke. Da Plato og Aristoteles havde ført
den græske Tænkning til dens høieste Udvikling, blev Dualismen
synlig, og Skjønhedsanskuelsens Opløsningstid begynder.
Bevidstheden om Spliden fører da først til Forsøg
paa at see Tilværelsens Heelhed i den ene Modsætnings Lys,
og idet de mislykkes, fremtvinger den dualistiske Theorier, i
% emphasis: „reent aandelig“ letterspaced (gaps 10--12 px in „reent“, 6--11 px in „aandelig“)
hvilke en \emph{reent aandelig} Gud faaer til sin oprindelige
Modsætning den dunkle Materie. Men den dualistiske Sondring
mellem Aand og Natur, der sønderrev Skjønhedseenheden,
maatte bidrage til at fremkalde Skepsis; den maatte
give Aanden, der ved en gabende Kløft var skilt fra det
Materielle, Mistillid til sin Erkjenden, til sin theoretiske Evne.
Og idet denne Tvivl vaktes, medførte den afgjørende Konsekvenser.
Medens før Aanden, trygt stolende paa sin Tankes
Kraft, havde sat den det høieste Maal og havde fundet Hvile
i dens Erkjendelse af det Guddommelige og af Tilværelsen,
søger den nu, i Mistillid til Tanken, andre Kilder for
Sandhed, og den troer at finde dem i det skjulte Urvæsens
Aabenbaringer gjennem de positive Religioner eller gjennem
dæmoniske Mellemvæsener, der forbinde den utilgængelige
Guddom med den endelige Tilværelse. Og med Mistilliden
til Viden maatte ogsaa følge Mistillid til Menneskets sædelige
Kraft. For Grækeren var Viden Menneskets Væsensbestemmelse;
derfor var hiin dobbelte Mistillid uadskillelig.
Der tyedes da atter til det Guddommelige og til de Mellemled,
der forbandt det med Mennesket, for at finde Bistand til
det, Aanden ikke ved egen Kraft kunde naae. Ogsaa her
% p. 444: ERRATUM.  The volume's own „Tilføielser og Rettelser“ (p. 572, volume PDF 578) reads (checked at the image): „S. 444, Linie 4 franeden staaer der en meningsforstyrrende Trykfeil, „Viden“ istedetfor „Veien“.“  The print reads „Viden“, transcribed as printed; the sense requires „Veien“.
% p. 444: this copy carries a reader's ink annotation keyed to that erratum: „Viden“ underlined, and in the left margin „Veien (s. 572)“; the manuscript marks are not transcribed.
bøiede Viden af fra Filosofien til Religionen, og den førte
gjennem Luttringer og Renselser, gjennem Flugt fra det
Sandselige til det guddommelige Væsen, fra hvilket al Bestemmelse
af Materie var udelukket. Tilværelsens Sandhed
% --- p. 445 ---
\opage{445}laae saaledes i det reent Aandelige, der hævede sig i uendelig
Afstand op over Materiens Ufuldkommenhed, i hvilken det
Ondes Spire søgtes, og Menneskeaandens Stræben, der i
Grækenlands klassiske Tid var tilfredsstillet indenfor det
Skjønhedshele, som gjenspeilede Ideen, gik nu udover Verden
til et Hiinsides, hvor Alt var Aand.

Med denne Udvikling i Aandslivet fulgte saa en Konsekvens,
der blev af væsenlig Betydning for Eftertiden. Idet
alt det Sandselige, Naturlige, tabte sin Betydning, tabte
ogsaa de Naturskranker deres Betydning, der adskilte Folkene;
den partikularistiske Aand, der skilte Grækere fra Barbarer, og
som kun i den fribaarne Græker saae den sande Repræsentant
for det Menneskelige, afløstes af en Universalisme, der anerkjendte
Mennesket \emph{som} Menneske, hævdede det Menneskelige
gjennem Menneske-Ligheden.

Hvad der ad denne Vei naaedes, fik Styrke ved de ydre
politiske Forhold. Det romerske Verdensrige gjennembrød
Nationaliteternes Skranker, knuste Folkenes Eiendommelighed.
Gjennem Lidelsen førtes Aanderne bort fra den utilfredsstillende
Endelighed, til det Hiinsidige, til hvilket Religionerne
viste hen. Men Folkereligionerne og deres individuelt udprægede
Gudeskikkelser vare allerede ifærd med at smelte
sammen til en uklar Eenhed, og denne Sammensmeltningsproces
begunstigedes ved den politiske Eenhed. Men medens
et Pantheon samlede de mange Guder og Theologien ordnede
deres Skarer, hævede der sig op over dem en høieste Eenhed,
løftet op over al Væren og al Tanke og dog seet som al
Værens og al Tankes Grund, og for dette usynlige, reent
aandelige Gudevæsen bleve de mange Guder tjenende Dæmoner,
der førte den higende Menneskeaand til dens Kilde, i hvilken
den sænkede sig ned.

Paa det Punkt, til hvilket vi have fulgt den græske
Aand i dens Udvikling, hvor den havde ombyttet Skjønhedsharmonien
med Spiritualismen, den nationale Partikularisme
med Universalismen, mødes den med den \emph{jødiske} Aand.
Denne havde i den gamle Hebraismes Tid trygt hvilet i det
% --- p. 446 ---
\opage{446}Nærværende og Dennesidige. Det udvalgte Folks Kontraktforhold
til Jehova, Loven med dens Belønninger og Straffe,
Offertjenesten og Kultus viste hen til et barnligt Udviklingstrin,
hvor Aanden har sin Norm udenfor sig og sit Maal
i et Endeligt. Ved Profeterne blev det Religiøse inderliggjort;
Tidens Trængsler, som de fortolkede fra et religiøst
Synspunkt, førte ud over det Nærværende til et Fremtidshaab
om et Messias-Rige; de første Anelser om en Udødelighed
glimtede frem, og Nationalbevidsthedens snevre Skranker
begyndte at aabne sig. Efter Exilet finde vi saa i en vis
Forstand en Tilbagegang, en ufri Fastklamren ved det Traditionelle.
Men paa den anden Side finde vi Forberedelser
til en fremskridende Udvikling. Vi finde saadanne i den
under persisk Paavirkning uddannede Forestilling om et
Aanderige, der er det jordiskes Forbillede og Sandhed, i Forestillingen
om en overnaturlig Messias og i Opstandelseslæren.
Endelig, i Tiden umiddelbart før Kristendommens Fremkomst,
see vi det Forberedte træde frem. Under Berøring med det
Græske udvikler sig nu bestemtere den dualistiske Sondring
mellem et Verdens og Materiens Rige, de onde Magters
Rige, og et længselsfuldt forventet himmelsk Rige i en naturløs
Virkelighed. Mod det himmelske Rige stræbes der hen
gjennem en asketisk Flugt fra Sandseligheden, og Udødelighedstroen
antager Formen af et Haab om en fuldkommen
Befrielse fra det jordiske Stof. Hermed staaer i Forbindelse
en sædelig Inderlighed og en almeen Menneskekjærlighed, der
praktisk førte til Fordømmelse af Slaveriet (Essener og Therapeuter).
Messiasforestillingen udvikles til det Overnaturlige.

I Opfattelsen af det Guddommelige paa abstrakt spiritualistisk
Maade; i den dualistiske Modsætning mellem det
aandeligt Guddommelige og det Materielle; i Menneskeaandens
dybe Trang til at finde Fred og Hvile i et Rige, der ikke
er af denne Verden, ved en Afdøen fra alt det Naturlige,
ved Flugt fra Verden; i Mistilliden til ved egen Kraft at
kunne naae dette Maal --- mødes saaledes ved den kristelige
Tids Begyndelse de to Folkeaander: den Aand, hvis Grund%
% --- p. 447 ---
\opage{447}tanke var \emph{Gud}, og den, hvis Grundtanke var \emph{Mennesket},
og Kristendommen, der forbinder de to Tanker i \emph{Gudmenneskets}
Tanke, bliver i sit Grundpræg bestemt ved sine
to væsenlige Faktorer.

Før vi paavise dette i Kristendommens Hovedlærdomme,
ville vi i Korthed see de \emph{positive Religioners begrebsmæssige
Udvikling} føre hen til det Punkt, hvor de to
Folkeaanders konvergerende Udviklingslinier have viist sig at
mødes.

Idet Menneskeaanden i den positive Religion, dreven af
sit Væsens Trang, søger et Forhold til det Uendelige, bliver
denne dens Uendelighedsstræben kun virkeliggjort under Endelighedsvilkaar:
gjennem en fremskridende Række af Former
maa den religiøse Bevidsthed nærme sig til sit Ideal. Det
Uendelige, den søger, finder den først i Naturen: Naturreligionen
er den positive Religions første Form, Naturen
det første guddommelige Væsen. Men den søgende Aand, der
endnu ikke formaaer at adskille sig selv som et Høiere fra
Naturen, formaaer ikke heller at adskille Naturen fra sig, saaledes
som den umiddelbart fornemmer sig selv som levende,
følende og villende. Naturen, for hvilken Mennesket bøier
sig i Tilbedelse, bliver ikke opfattet \emph{som} Natur, men bliver
menneskeliggjort, bliver opfattet som et personligt, levende,
menneskelignende Væsen. Saaledes seer Barnet de ydre Ting
som levende og besjælede; saaledes kjærtegner det dem eller
vredes paa dem, fordi det tillægger dem Villie og Tanke.
Idet Naturvæsenet, der tilbedes, menneskeliggjøres, bliver
det ved selve denne Personifikation i en vis Forstand et
% p. 447: only „over“ of „overnaturligt“ is letterspaced (gaps o-v 12, v-e 11, e-r 8 px; r-n 4 px solid)
\emph{over}naturligt Væsen. Naturgjenstanden bliver ligesom fordoblet:
over den sandselige Naturting hæver sig Tingens
Genius, dens Fylgie, dens Ferver, --- for at minde om de
forskjellige Religioners Udtryk. Som Dryaden skiller sig fra
Træet og dog er Træet, som Najaden skiller sig fra Strømmen
og dog er Strømmen, saaledes foregaaer overalt Fordoblingen i
det Enkelte og i det Omfattende. Hvad der i Fordoblingen staaer
ligesom bagved det Naturlige, bestemmes saa som et Usyn%
% --- p. 448 ---
\opage{448}ligt, som et for Sandserne Utilgængeligt, --- saaledes undgaaer
Naturreligionen den Modsigelse, der i den er mellem
Forestilling og Virkelighed, Indbildning og Sandhed. Modsigelsen
mellem de to Objekter, i hvilke det ene Naturobjekt
har skilt sig, skal herved hæves, idet det Objekt, der er Religionens
egenlige Gjenstand, unddrages fra det umiddelbare
Forhold, bliver et Væsen, der kun er Gjenstand for Troen,
for Aanden. Herved er allerede Aandsreligionens Opfattelse
af det Guddommelige antydet. Og Aandsreligionen afløser
Naturreligionen, hvor den menneskelige Udvikling er skredet
saavidt frem, at Mennesket fra et blot naturbestemt Væsen
er udviklet til et Væsen, der adskiller sig fra Naturen og
ligesom koncentrerer sig i sig selv. Naar Mennesket, gjennem
det menneskelige Samfundsliv, bliver til et Væsen, der bestemmes
ved bevidste Fornuftlove, naar andre Magter end
Naturmagter: politiske, moralske Magter, blive de væsenlige
Gjenstande for dets Afhængighedsfølelse, da viser Verden og
Naturen sig ikke længer for det som Gud, men som et ved
en guddommelig Forstand og en guddommelig Villie bestemt
og af den afhængigt Væsen. Idet Mennesket selv ved Villie
og Forstand hæver sig over Naturen og saaledes bliver et
supranaturalistisk Væsen, bliver med det Samme ogsaa dets
Gud et supranaturalistisk Væsen, ikke længer Eet med Naturen,
men Herre over Naturen; og han maa endelig, for
i streng Forstand at være Herre, blive Skaber af Naturen.
Og de mange Guder, som Aandsreligionerne modtage fra
Naturreligionerne, maae konsekvent smelte sammen til en guddommelig
Eenhed. Den menneskelige Bevidstheds og Aands
Eenhed, der omfatter hele den ydre Verdens Mangfoldighed,
er det, der fører til at statuere Guds Eenhed.

Saaledes arbeider den positive Religion sig frem fra
Naturreligionernes laveste Trin til Aandsreligionernes høieste,
til Forestillingen om den ene, aandelige, skabende Gud. Men
netop naar den naaer til dette Punkt, bliver det Guddommelige
opfattet abstrakt spiritualistisk som det \emph{naturløst Aandelige}.
Den positive Religions Organ er ikke den begri%
% --- p. 449 ---
\opage{449}bende Tanke, men Forestillingen og Følelsen. Hvad den besidder
sammentrængt i Følelsens rige Inderlighed, giver den
Udfoldelse i Forestillingen; kun i denne Udfoldelse faaer
Følelsesindholdet Bestemthed og Klarhed. Men Forestillingens
% p. 449: „sandelige“ as printed (for „sandselige“?)
Særkjende er det, at den fra den sandelige Anskuelse, hvilken
den skal føre over til Tanken, har bevaret en Udvorteshedens
Form, der bevirker, at dens Indhold træder ud i en endelig
Sondring og for det indre Øie fæstnes i denne Sondring.
Derfor bliver det Guddommelige, al Tilværelses Kilde og
Grund, stillet som guddommelig Personlighed over og udenfor
% p. 449: again only „over“ of „overnaturligt“ letterspaced (gaps 11, 12, 8 px; r-n 9, then 4--6 px)
denne Tilværelse, bestemt som \emph{over}naturligt Væsen, i den
Betydning, at dette Overnaturlige danner \emph{Modsætningen}
til Naturens positive Bestemmelser. Det Guddommeliges
\emph{Forestillingssondring} fra Naturen som det Overnaturlige
indeslutter Opfattelsen af det som det naturløse Aandelige:
Forestillingstranscendensen indeslutter den abstrakte Spiritualisme.
Men naar Opfattelsen af det Guddommelige som
den ene, aandelige, skabende Gud er de positive Religioners
Maal, saa maa altsaa den abstrakte Spiritualisme være
Kjendemærket for den høieste af de positive Religioner,
og den høieste og afsluttende af dem er Kristendommen.
Kun ved saa konsekvent, som det overhovedet lader sig gjøre,
at gjennemføre en abstrakt spiritualistisk Opfattelse, faaer
Kristendommen en eiendommelig Livsanskuelse. Alle forudliggende
Stadier, alle Forberedelser ere allerede givne i de
for-kristelige Religioner; kun hiint bliver tilovers for den.

I de kristelige Grundforestillinger er overalt det abstrakt
Spiritualistiske let at paavise. Det træder frem i Læren
om \emph{Gud} som den rene Aand, der boer i det Lys, i hvilket
intet Materiens Mørke er. --- Det træder frem i Forestillingen
om en \emph{Skabelse af Intet} ved den almægtige Villie:
thi derved sættes det Skabte som det, der i sit Væsen er et
Intet, som det, der kun fra Almagtens Vilkaarlighed, ved
hvilken det bæres oppe, har en laant Væren, der hvert
Øieblik kan tilintetgjøres, som det partielt skeer i Underet, for
hele Naturen i den Verdensundergang, i hvilken det, der
% --- p. 450 ---
\opage{450}var et Intet, vender tilbage til det Intet, af hvilket det er
kommet, og kun lader det Aandelige bestaae. --- Det træder
frem i hele Forestillingen om det Guddommeliges Virken
som en \emph{undergjørende} Virken, thi i Underet virker Almagten
uden alt Realt, benægter den det Reales Væren og dets
Gyldighed. --- Det træder frem i Læren om den \emph{menneskelige
Livsopgave}, der bliver den, at søge Eenheden med Gud
gjennem Kristus ved Afdøetheden fra Verden, ved Fornægtelsen
af den Natur, der „ligger i det Onde“, hvilken, naar
den gjør sig gjældende som Natur i os, bestemmes som det Syndige,
og som derfor bliver det, til hvilket Aanden kun skal forholde
sig negativt, saaledes at det ikke bliver et Stof, der
skal formes, men et saadant, der skal bortkastes, idet den
Kristne skal flye det naturlige Livs Forhold, som for ham
ingen Betydning have, men kun kunne blive besmittende,
fordi de ikke kunne sættes i et positivt Forhold til det abstrakt
Aandelige (saaledes Ægtestand, Statsliv, Eiendom). ---
Det træder frem i den til denne Opfattelse af Livsopgaven
svarende Opfattelse af \emph{Livet efter Døden}. Hvad der for
Bevidstheden staaer som et Ideal, men som et Ideal, den
mistvivler om at kunne naae i denne Tilværelse, det lægger
den hen til et hiinsidigt Liv, hvor alt det Naturlige, der
endnu hæmmer og tynger det Aandelige, skal være forsvundet,
hvor Aanden skal være frigjort fra enhver Naturbegrændsning,
og hvor selv det Legeme, i hvilket Fantasien iklæder den
henfarne Aand for at kunne anskue den, skal være et „aandeligt
Legeme“. --- Det træder endelig frem i selve Læren om \emph{Gudmennesket},
netop i det Centralpunkt i Kristendommen,
hvor det Guddommelige forbindes med det Menneskelige og
gjennem det Menneskelige synes ogsaa at træde i et positivt
Forhold til det Naturlige. Thi dette Skin forsvinder, saasnart
der sees hen til, at Gudmenneskets Liv heelt igjennem
er et \emph{Under} og ved Underet er løsnet fra det menneskelige
Livs sædvanlige Natur-Betingelser: ved Fødselen og i Døden,
og paa hvert Punkt mellem disse. Eenheden af Gud og
% --- p. 451 ---
\opage{451}Menneske i Gudmennesket bliver kun Eenheden af Gud og
det \emph{naturløse} Menneskelige. ---

Saaledes ere vi, ad forskjellige Veie, komne til det
samme Resultat: at Kristendommens Karakteer er den naturløse
Aandelighed, der udpræger dens Opfattelse af det Guddommelige
og af det Menneskelige. Betragtningen af de
historiske Forudsætninger, af de positive Religioners begrebsmæssige
Udvikling og af selve Kristendommens Grundlærdomme
have ført os til een og samme Opfattelse og ved
denne Samstemmen givet den Sikkerhed. Og de svage Forsøg,
der fra Mæglingstheoriens Side gjøres paa at bestride denne
Opfattelse, ere af den Art, at de ikke let ville kunne gjøre Nogen
tvivlraadig. Vi ville i Forbigaaende --- thi det vil være
tilstrækkeligt --- kaste et Blik paa et saadant Forsøg, der
findes i den Samling af opbyggelige Betragtninger over ethiske
Emner, som Biskop Martensen nylig har udgivet.

Der begyndes med det Postulat, at den hellige Skrift paa
ethvert af sine Blade skal bevidne, at Gud ikke er den naturløse
Aand. „Skriften kjender kun den levende Gud. Men den
levende Gud kunne vi (!) ikke tænke anderledes end i Forhold
til en evig, ham underordnet Natur“. Der fortsættes
saa: „Naar Skriften taler om Gud som et \emph{uopløseligt}
Liv (Hebr. 7) saa kunne vi heller ikke tænke os dette uden
som den uopløselige Eenhed af Livets Modsætninger, som
hvilke vi kun kjende Modsætningen mellem Aand og Natur,
det Ethiske og det Fysiske“. Denne Exegese, der ved sin
Blanding af Barnagtighed og Spidsfindighed minder om
Rabbinernes, er saa uheldig som mulig, hvilket Enhver vil
kunne overbevise sig om ved at gjennemlæse to Linier paa
% p. 451: no full stop after „(V. 16 f.)“ before „Kristus“, as printed
det citerede Sted (V. 16 f.) Kristus bliver sammenlignet
med Melchisedek, som den evige Ypperstepræst i Modsætning
til den dødelige Ypperstepræst efter Mose Lov. Som den,
der ikke har „Dagenes Begyndelse eller Livets Ende“, som
den, der „bestandig bliver Ypperstepræst“, kaldes han Ypperstepræst
„efter et uopløseligt Livs Kraft“, og dette Udtryk fortolkes
i det umiddelbart Efterfølgende, idet der siges: „\emph{thi}
% --- p. 452 ---
\opage{452}Du er Præst til \emph{evig} Tid efter Melchisedeks Orden“. Her
er altsaa ganske simpelt en Modsætning mellem det Liv, der
kan ophøre i Døden, og det, over hvilket Døden ikke har
nogen Magt; men ikke en eneste Vending tyder hen til hiin
høiere Eenhed af „Livets Modsætninger“, som vi „maae
tænke“. Man maa være ualmindelig blottet for Grunde for
at tye til en saadan Fortolkningskunst. Dog, der er endnu
et Argument tilbage, der næsten overgaaer de forrige. Det
% p. 452: „Anthropomofismer“, „Anthropomofismes“, „Anthropomofistiske“ (twice) all as printed, without r
hentes derfra, at Skriften overalt taler om Gud i Anthropomofismer.
„Men den religiøse Anthropomofismes Sandhed
beroer paa Naturen i Gud“. „Den levende Fromhed lader
sig det ikke afdisputere, at Gud har Øine og Øren, Hænder
og Fødder, en Arm, der ikke bliver forkortet; og den opgiver
ikke at see hans Finger i Verdensbegivenhederne og i den
Enkeltes Liv; thi om den end er sig det Billedlige, det Symbolske
bevidst og erkjender, at Alt maa fjernes, hvad der
kun hører til den skabte Indskrænkning, dog vedbliver den at
fastholde, at der til Alt dette maa være et virkeligt Tilsvarende
i Gud; hvilket med andre Ord vil sige, at Gud maa
have Aabenbaringsorganer, Redskaber --- en Organisme“. ---
Altsaa: det Anthropomofistiske skal være Billede, Symbol,
altsaa \emph{ikke} udtrykke Virkeligheden; men dog skulle vi fordre
et \emph{virkeligt} Tilsvarende, og medens vi fjerne Alt, der kunde
give disse Bestemmelser Indhold, skulle vi tillægge Gud et
Legeme, en Organisme. Det er den gamle, ovenfor karakteriserede
Methode: der gives med den ene Haand og tages
med den anden. Det Anthropomofistiske skal være ikke-virkeligt
og dog virkeligt, uegenligt og dog egenligt, og det for
at støtte en stakkels vaklende, halvspekulativ Theori, der fabler
om „fysiske Almuligheder“ i Gud, om et „relativt Intet“,
der skal være ligt med de „evige, fysiske Muligheder“, om
% p. 452: only „over“ of „overnaturlige“ letterspaced
„fysiske, men \emph{over}naturlige Kræfter“, om en Materie, der
er forskjellig fra Naturen ved at være „udtraadt (hvorledes?)
af Forbindelsen med Aanden“, o. desl. Det er karakteristisk,
at Theoriens Forkynder selv maa indrømme, at der „ingen
konkret Begriben“ kan finde Sted af den. Tilvisse! de
% --- p. 453 ---
\opage{453}tomme Fraser opløse sig af sig selv i et Intet, der ikke er
et „relativt“, men aldeles fuldstændig et „absolut“, og det er
ved en „uforsvarlig Misforstaaelse“, at den hellige Skrifts
Forfattere faae Skyld for at vedkjende sig en saadan Theori.
Hvad vilde vel Johannes-Evangeliets Forfatter sige til den
anførte Fortolkning af „Skriftens Billedsprog?“

% p. 453: short centred rule (c. 1.8 cm) divides the text
\begin{center}\rule{1.8cm}{0.4pt}\end{center}

Naar vi fastholde Opfattelsen af Kristendommen som
den, hvis Karakteermærke er den naturløse Aandelighed, saa
bliver det historisk givne Fænomen, der træder os imøde
med Kristendommens Navn, os let forklarligt. Vi forstaae
ikke blot Kristendommen i dens Kamp med Verden, i dens
klassiske Skikkelser, i de Livsformer, i hvilke den skarpt har
udpræget sit eiendommelige Væsen; men vi forstaae ogsaa de
Former, hvor det Kristelige er blandet med sin Modsætning.
Vi forstaae uden mindste Vanskelighed, hvorledes den naturfornægtende
og verdensfornægtende Religion, --- idet den
(tvertimod den første kristne Tids Haab om Kristi nære
Gjenkomst) skulde vedblive at existere i den fremmede Verden,
idet den skulde forkyndes i denne Verden, idet den skulde
værge sig mod Angrebene fra den, og idet den skulde udfoldes
i Bevidstheden --- maatte ombytte sit oprindeligt kun
negative Forhold til Naturen, til de naturlige Evner og
Kræfter (deri den naturbestemte Tænkning indbefattet) og til
Verdenslivets konkrete Former med en nødtvungen Forbindelse
med det Fornægtede, der optoges som Midler for det
Kristelige. Saaledes forklare vi det, man har kaldt en
kristelig Videnskab, en kristelig Kunst, en kristelig Sædelighed,
en kristelig Stat. Men vi forstaae tillige, at den
nødtvungne faktiske Anerkjendelse af det i Principet Fornægtede
aldrig kunde blive en fuld og sand Anerkjendelse, og at
derfor den nødtvungne Forbindelse ikke kunde blive en virkelig
Forening; thi det Optagne blev, som blot Middel, uden
Væsenssammenhæng med det Kristelige. Saaledes viser ogsaa
% --- p. 454 ---
\opage{454}Forholdet sig overalt, hvor det Kristelige bliver sig selv
bevidst, saaledes som vi i det Foregaaende (H. 4, S. 291)
have paaviist med Hensyn til den „kristelige“ Filosofi.

Men medens Fænomenet fra dette Synspunkt, der
sætter Modsætningen mellem det Kristelige og det Humane,
bliver let forklarligt, støder derimod den Mægling, den
spekulative Theori forsøger paa Grundlag af den ovenfor
karakteriserede Opfattelse af Kristendommens Forhold til det
Naturlige, paa uløselige Vanskeligheder. Gjøres der virkelig
Alvor med Mæglingen, saa bliver Kristendommen, der
klart og tydelig forkynder sig selv som en ny Skabelse, d.
v. s. som en absolut ny Begyndelse, og som i Aarhundreder
er opfattet som en saadan af sine Bekjendere, et reent Overflødigt,
da det, den skal være i Eenheden med det Humane,
allerede har været givet forud for Kristendommen. Og alle
de store klassiske Skikkelser gjennem Kristendommens hele
Tid maae saa opfattes som eensidige, begrændsede Individualiteter,
for hvilke Kristendommens egenlige Sandhed ikke er
gaaet op. --- Gjøres der Indskrænkninger og tages der Forbehold,
saa at Mæglingen kun bliver halv --- og kun saavidt
vover i Regelen den spekulative Theologi at gaae, naar
den skal udtale sig om Principerne, --- saa forkvakles baade
det Humane og det Kristelige, og begge blive meningsløse.
Det Humane forkvakles og bliver meningsløst derved, at
det, trods alle anerkjendende Fraser, kun faaer en Skinvirkelighed
og en Skinbetydning. Det Kristelige forkvakles
derved, at de Elementer, det faaer i sig fra det Fremmede,
skjøndt de ligesom øieblikkelig sublimeres, dog kaste et Skjær
over det Kristelige, der berøver det dets specifiske Eiendommelighed,
gjøre det til et Principløst og derfor Meningsløst.

% p. 454: short centred rule (c. 1.7 cm) divides the text
\begin{center}\rule{1.8cm}{0.4pt}\end{center}

Men staaer da Opfattelsen af det Kristelige som den
naturløse Aandelighed fast, saa staaer det ogsaa fast, at der
maa vedblive at være en Modsætning mellem det og det
% --- p. 455 ---
\opage{455}Humane, trods den Betydning, det Kristelige kan have havt for
Udviklingen og Luttringen af dette. Kristendommens Tid har
været Menneskeaandens Skjærsild, det Humanes Luttringstid.
Den naturløse, abstrakte Aandelighed med sin negative Uendelighed
har ligesom udbrændt Sandseligheden og Endeligheden
af Menneskevæsenet og derved gjort det muligt for Aanden,
idet den fordybede sig i sit Væsens Rigdom, at give denne
en indholdsrig Uendeligheds Form og saaledes at kunne leve
et evigt Liv i Aandens Forsoning indenfor Virkeligheden.
Den har kunnet nedbryde Skrankerne, der skilte Folkene fra
hverandre, og styrke den Bevidsthed om en Menneskehed, til
hvilken allerede Oldtiden var naaet. Og Kristendommen
har ved kraftigt at minde Menneskeaanden om, at Menneskelivet
først finder sin Sandhed i Forholdet til Gud, givet
den Impulsen til ligesom at fordybe alle det sædelige Livs
humane Forhold i Gudsforholdet, at gjøre dem alle til religiøse
Forhold. Men den naturløse Aandeligheds Livsanskuelse
indeslutter Konsekvenser, mod hvilke den humane Livsanskuelse
maa værge sig. Idet Menneskelivets Naturgrundlag
bliver et Intet, og nærmere et saadant Intet, der, naar det
gjør sig gjældende som et Noget, som en Virkelighed, sees
som det Onde, saa er Grundlaget for det konkrete sædelige
Liv, for alle de Livsformer, der, som Ægteskabet, Samfundet,
Staten, hvile paa en Naturbasis, borttaget: kun
Askesens Flugt fra Verden, kun en Forsagen af alt det sædelige
Livs Virkelighedsindhold bliver tilbage. Og naar
Naturen forsvinder fra den kun aandelige Guds Væsen, saa
bliver Guds Virken i den ved en ubegribelig Almagtens
Akt satte Natur den undergjørende Virken, der ikke levner en
Plads for Naturlove, og med Naturens Love falde Tankens
Love, og Tankens Gyldighed forsvinder. Overalt møder
Tanken Underet som den høiere Sandhed, og i Kristendommens
Centralpunkt: i Læren om Gudmennesket, om det evige,
altomfattende Guddommeliges Eenhed med det enkelte, under
Endelighedsbetingelsernes Begrændsning existerende Menneske,
møder den Underet i dets prægnanteste Form, og den kan
% --- p. 456 ---
\opage{456}kun hævde sig selv og det menneskelige Væsen, hvis Grundbestemmelse
den er, ved at skille sig fra, hvad der maa tilintetgjøre
den. ---

Men idet Sondringen saaledes træder frem mellem den
humane Bevidsthed og den Religion, der som de historisk
givne positive Religioners høieste og afsluttende Form er
den positive Religions berettigede Repræsentant, vender den
humane Bevidsthed sig ikke bort fra det Religiøse i dets ideale,
for det Positives Begrændsning frigjorte Form, men den
finder dette ideale Religiøse, idet den fordyber sig i sig selv,
idet den trænger ind til det menneskelige Væsens Grund, til
det Punkt, hvor det er knyttet uopløseligt til sit guddommelige
Ophav. Hvorledes dette Forhold mellem det Humane
og det ideelt Religiøse nærmere bliver at bestemme, skulle vi
søge at vise i et følgende Afsnit, der vil afslutte disse Udviklinger.

% p. 456: signature set right in bold Fraktur; then, about 2.7 cm lower, a single centred rule (c. 2.6 cm), not a double rule, closes installment II; the rest of the page is blank but for show-through
\begin{flushright}\textbf{Constans.}\end{flushright}
\begin{center}\rule{2.6cm}{0.4pt}\end{center}

\clearpage
% ===== III. (pp. 492--510) =====
% --- p. 492 ---
% p. 492: upper page blank (no tail of a preceding article); head printed over two lines in display blackletter ("Det Kristelige og det Humane som Modsætninger" / "i vor Tids Bevidsthed."), short rule beneath, then the centred numeral
\opage{492}\begin{center}\textbf{\large Det Kristelige og det Humane som Modsætninger i vor Tids Bevidsthed.}\end{center}

\secnum{III.}

Naar vi i de mange Former, under hvilke den humane Livsanskuelse
--- i dens klarere eller dunklere bevidste Modsætning
til den kristelige --- fra den nyere Tids Begyndelse har givet
sig et Udtryk, ville søge det Fælles og Gjennemgaaende, der
maa sættes som det væsenlig Bestemmende, da ville vi finde
dette i to Hovedpunkter. Det ene er den humane Livsretnings
Stræben efter at hævde den \emph{uforstyrrede, harmoniske
Eenhed i det menneskelige Væsen}; det andet er dens
Stræben efter at hævde \emph{det Ethiske i dets Konkretion},
at give det sædelige Liv et Virkelighedsindhold, rigt nok til
at udfylde \emph{hele} Menneskelivet. Men disse to Hovedpunkter
staae, som det let vil vise sig, i den nøieste indre Sammenhæng.

Medens den kristelige Livsanskuelse --- ved at fornægte det
Naturliges Gyldighed og Væsenlighed og ved at sætte Sandheden
i det, der for Tanken var det Utilegnelige, det, der
stred mod alle Erkjendelsens Love --- havde sat et Brud i
Menneskevæsenet, et Brud, der blev uendelig skærpet, idet
Læren om Arvesynden satte Menneskenaturen som forstyrret og
fordærvet: hævder den humane Livsanskuelse Eenheden i den
\emph{dobbelte} Betydning af en Eenhed mellem Natursiden og
% --- p. 493 ---
\opage{493}Aandssiden i Mennesket og af en Eenhed, en Harmoni
mellem Viden og Tro, mellem Erkjendelse og Villie, mellem
det Theoretiske og det Praktiske. Den Natur, som den middelalderlige
Aand havde fornægtet, som den ængstelig gyste tilbage
for, medens Fantasien uimodstaaelig lokkedes af det
Ængstende, \emph{den} blev ved den nye Tids Begyndelse atter
indsat i sin Ret. Idet Menneskeaanden søgte ind i sig selv
for at finde sig selv i sin Frihed og Selvstændighed, kunde
den ikke mere blive blind for det uløselige Baand, der bandt
den til Naturen, for den Livsfylde, den hentede fra denne.
Den angstfulde Svæven mellem Frygt og Tillokkelse blev
afløst af begeistret Hengivelse. Det var, som om Guddommen
atter var vendt tilbage til den forladte Natur, der nu atter
saaes i Væsenssammenhæng med sit Ophav, atter i sin
Skjønhed blev Gjenstand for Aandens Beundring, atter i
sin Væsenlighed og Lovbegrundethed blev en Erkjendelsens
Kilde for Aanden. Denne Retning mod Naturen gaaer
gjennem hele den nyere Tid. Vi finde den i Renaissancetidens
begeistrede Forguden af det naturlige Universum; vi
finde den i Erfaringsvidenskabens Hengivelse til Naturen; vi
finde den i det 18de Aarhundredes Natursværmeri og i dets
Fordring om Naturlighed i Opdragelsen og i Livet; vi finde
den i vort Aarhundrede i Fornyelsen og Udbredelsen af
Materialismen, den theoretiske og praktiske, i Naturvidenskabernes
mægtige Udvikling, i Filosofiens Bestræbelse for at
see det ubevidste og bevidste Sjæleliv i Mennesket som en
Eenhed og for at give det Guddommelige et Naturindhold, ja
vi finde den selv i den filosoferende Theologis famlende
Forsøg paa at faae et Fantasibillede af Naturen ind i den
overnaturlige Gud. Men den i sin Væsenlighed atter
anerkjendte Natur bliver saa, som Naturbestemthed i det
Menneskelige, seet i Eenhed med det Aandelige. Det organiske
Liv udfolder sig til et aandeligt Liv, til et Bevidsthedsliv;
fra vor Naturside, gjennem de organiske Betingelser,
faaer vor Erkjendelse sit Indhold; fra vort Driftliv
faaer Villieslivet sit Stof. Det Naturlige i os bliver ikke
% --- p. 494 ---
\opage{494}længer det, der kun skal fornægtes, det, hvorfra vor Villie
skal afdøe for at leve et reent aandeligt Liv; det bliver atter,
som i den græske Oldtid, det, der skal formes og harmoniseres
af den Villie, der ved Erkjendelsen af vort Væsen og dets
Grundforhold har faaet Klarhed over sin Opgave, og som
virkeliggjør den sande Erkjendelses Fordring. I eensidige
Former træder denne Anerkjendelse af det Naturlige i os frem
i de psykologiske, ethiske og sociale Theorier, der lade Driftbestemtheden
i os blive det høieste Raadende, der lade den
beherske vor individuelle Udvikling og Samfundsforholdenes
Ordning. Men hvad der her paa en eensidig Maade gjør
sig gjældende, faaer sin sande Form, hvor en oprindelig
Lovmæssighed i vort \emph{aandelige} Væsen anerkjendes, efter
hvilken Naturbestemmelserne i os ordnes og formes, under
hvilken de tjenende bøie sig, idet de anerkjendes i deres
Berettigelse.

Idet Eenheden hævdes af \emph{Natursiden} og \emph{Aandssiden}
i os, hævdes tillige Eenheden mellem \emph{Viden} og \emph{Tro},
mellem Erkjendelse og Villie. Det er Fornægtelsen af Naturen,
der berøver Tanken dens Fodfæste; det er Underet, i hvilket
det Reales Virkelighed er tilintetgjort til Fordeel for den
overnaturlige Almagts ubegribelige Virken, som sætter Tanken
en uoverstigelig Skranke og berøver dens Love deres Sandhed.
Hvor Troen, som i den positive Religion, er Tro paa den
undergjørende Gud, paa den Gud, for hvis Villie Naturens
Love forsvinde, paa hvis Bud Naturen selv forsvinder, ligesom
den er bleven til paa det: \emph{der} maa Tro og Viden danne
uforenelige Modsætninger, og idet Troen som den høiere
Sandheds Organ dømmer Viden til Underkastelse, til Selvopgivelse,
fornægtes her en Grundform af Menneskeaanden,
en Væsensbestemmelse i den. Der kommer da et Brud i
Væsenet, en indre Splid; thi det Væsenlige, der skal fornægtes,
kan, netop som væsenligt, ikke tilintetgjøres. Anderledes
kan Forholdet blive, naar det Naturlige hævdes i sin
Sandhed, naar Naturen ikke længer ved en oversvævende
Almagt holdes oppe over Intets Afgrund, men naar den,
% --- p. 495 ---
\opage{495}som bærende det Guddommeliges Liv i sig, hviler i sine
væsensbestemte Love. Her bliver det Guddommeliges Væsen
og Virken ikke et Saadant, som Tanken kastes tilbage fra,
og som kun kan gribes af en Tanken fornægtende Tro: det
bliver et Indhold for den Viden, der, følgende den i Menneskevæsenet
liggende Fordring, fra det Endelige i dets Mangfoldighed
og Begrændsning tillidsfuldt vover sig ud til det
Uendelige. Og idet Viden og Tro ikke længer staae som
Modsætninger, idet Menneskeaanden finder Hvile i sin Erkjendelses
Indhold, bliver der ogsaa en Harmoni mellem Erkjendelse
og Handling. Den Tro, der stod imod Viden, gav
ogsaa Villien en Lov, modsat Videns Lov. Ikke i Menneskevæsenets
Fordring, men i en overnaturlig, ubegribelig Villies
% p. 495: a risen space (quad) prints as a faint mark before "hvor"; not a character, not transcribed
Bud søgtes Handlingens Norm, hvor Handlingens Opgave
ikke var den reent negative: at flye det Naturlige. Men naar
Tro og Viden ere i Harmoni, saa sees ogsaa Videns Fordring,
der udtrykker Menneskevæsenets Lov, Menneskevæsenets
egen Fordring, som den ene gyldige Fordring til Villien;
saa sees det samme, uadsplittede Menneskevæsen, der kommer
til Klarhed over sig i sin Erkjendelse, som udtrykkende sig, i
Overeensstemmelse med sig, i den bevidste Handling. Det
Theoretiske og det Praktiske sees da i Harmoni: eet Væsen
udfolder sig gjennem begge. Aanden vakler ikke mellem et
dobbelt Princip, en dobbelt Lov; den nages ikke af Tvivl
om sin Erkjendelses Sandhed og sin Handlings Gyldighed:
i Harmonien har den sin Betryggelse; den har Hvile i Eenheden
med sig selv. ---

Som det \emph{andet} Hovedmotiv i den humane Livsretning
bestemte vi \emph{Hævdelsen af det Ethiske i dets Konkretion},
% p. 495: only "Virkeligheds" is letterspaced; "indhold" is set solid
af det sædelige Livs \emph{Virkeligheds}indhold. Indenfor
den Religion, hvis Karakteer var den abstrakte, naturfornægtende
Aandelighed, kunde den ethiske Livsopgave konsekvent
kun blive en Afdøen fra alt det Naturlige, en Askesens Flugt
fra Verden for at finde Hvile i et reent aandeligt Liv. Idealet
maatte blive et hiinsidigt Liv, hvor alle Naturskranker vare
forsvundne, hvor alt det Tyngende og Besmittende i den
% --- p. 496 ---
\opage{496}naturbestemte jordiske Tilværelse var fjernet. Saalænge
Mennesket levede under det jordiske Livs Vilkaar, maatte det
bestandig i sin Bundethed til det Naturlige føle et hæmmende
Baand og en Spire til Ufred i sit Indre. Det Naturliges
Drifttilskyndelser, gjennem hvilke Naturen hævdede sig som
en Magt i os, føltes som Fristelse og Besmittelse, og
længselsfuldt vendte Blikket sig mod den Fremtid, i hvilken
Aanden skulde leve frigjort fra dette Hæmmende. Den kristelige
Oldtids og Middelalderens Askese, Klosterlivet og Eneboernes
Flugt fra Verden, var ikke Eensidigheder og Overdrivelser,
men konsekvente Udtryk for Religionens Aand. Hvor det
Naturlige har mistet sin Gyldighed, hvor den naturløse Aand
er det Høieste, \emph{der} kan der konsekvent ikke leves et andet Liv
end det, som i hvert Øieblik er en Afdøen fra Verden. For
de konkrete Former af det sædelige Liv kan der ikke blive
nogen Plads; deres Gyldighed er betinget ved Gyldigheden
af Naturgrundlaget i os. Alle de sædelige Forhold og alle
de sædelige Samfundsformer vise, direkte eller indirekte, tilbage
til dette Naturprincip; de falde, naar det fornægtes.
Men Fornægtelsen af det er for Mennesket en Selvmodsigelse,
fordi det menneskelige Væsen i uadskillelig Eenhed
omfatter det Aandelige og det Naturlige. Den Kraft, ved
hvilken Naturen i os skal kues ned, kan kun komme fra
selve Naturgrunden i os, der er blivende Bestanddeel i vort
Væsen: i selve Benægtelsen af det Naturlige bekræfter det
sig selv. Den naturløse Aandeligheds Liv bliver derfor bestandig
kun et Forsøg, der idelig kan gjentages, men aldrig kan
gjennemføres, fordi det, der skal dødes, i selve Kampen faaer
Kræfter fra den skjulte Grund. Derfor er der Angst og
Uro i dette Liv, og den kan kun fjernes, naar Alt, hvad
der hører med til Menneskevæsenet, anerkjendes, naar dette
efter sit hele Omfang bliver Gjenstand for den ethiske Handlen,
naar Opgaven bliver den frie Virkeliggjørelse af Menneskevæsenet
efter dets Fuldstændighed, naar \emph{hele} Livet i dets
Virkelighedsfylde drages ind under det Ethiske.

% --- p. 497 ---
\opage{497}En Stræben herefter finde vi ogsaa, i mange Former,
ufuldkomnere og fuldkomnere, gjennem hele den nyere Tid.
Det er en saadan Stræben, der træder frem i det borgerlige
Livs verdslige Sædelighed, som ved Middelalderens Slutning
emanciperer sig fra det Kirkelige; det er den, der træder frem,
naar Religionen forvandles til Moral og i denne Forvandling
tidt trivialiseres; det er den, der gaaer igjennem de filanthropiske
Bestræbelser, gjennem Socialismens Forsøg paa at
ordne Samfundsforholdene, gjennem Filosofiens Bestræbelser
for at arbeide Ethiken ind i Virkeligheden, for at bestemme
Principet for den sædelige Handlen saaledes, at det sigter
hen mod Livets hele Rigdom.

Men idet den humane Livsretnings Stræben efter at hævde
det sædelige Liv i dets konkrete Virkelighed viser tilbage til
Hævdelsen af det Naturliges Gyldighed, viser den derigjennem
tilbage til, hvad vi udhævede som det \emph{første} Motiv: Hævdelsen
af Væsenets Eenhed. Det er det Menneskevæsen, som i sin
udeelte Eenhed harmonisk forbinder Natur og Aand, der kan
udfolde sig i et konkret sædeligt Liv. Men Menneskevæsenets
\emph{Eenhed} er uadskillelig fra Menneskevæsenets \emph{Uforstyrrethed},
og idet den humane Livsretnings Stræben sigter hen mod
det dobbelte Maal, vi have angivet, og som ligger i samme
% emphasis: „det uforstyrreligt Gode i Menneskenaturen“ letterspaced (gaps 10--15 px; „Gode“ 10--12 px against 3--6 px in the solid „Troen paa“)
Linie, bæres den af Troen paa \emph{det uforstyrreligt Gode
i Menneskenaturen}, paa Menneskenaturens \emph{væsenlige
Fuldstændighed} under al den Ufuldkommenhed, der er
uadskillelig fra dens Udvikling og dens Kampe. Den faste
Overbeviisning, at Menneskevæsenet ikke ved sin Gjerning, der
som Udtryk af Væsenets Virksomhed er en Bekræftelse af
Væsenet, kan forvandle, negere dette Væsen, er \emph{gjennemgaaende}
Særkjende i alle Former af den humane Livsanskuelse.

\begin{center}\rule{1.8cm}{0.4pt}\end{center}

Idet den humane Livsretning principielt sondrer sig fra
den positive Religion, kan det ikke undre os, at den hyppigt
har ført til en Sondring fra det Religiøse overhovedet, til
% --- p. 498 ---
\opage{498}en „religionsløs Sædelighed“, ja endog til en lidenskabelig
Kamp mod Religionen. Dette er dog hverken et Gjennemgaaende
i den humane Livsretnings mange Former, ikke heller
er det den naturlige Konsekvens. I mange af hine Former,
f. Ex. i socialistiske Theorier, finde vi en, om end uklar,
Stræben efter at give Livsanskuelsen en religiøs Afslutning,
der her i Regelen ikke undgaaer det Fantastiske. I andre Former
finde vi ligesom en Rest af det positivt Religiøse, rationaliserede
Fragmenter af dette. I nogle endelig, og disse betragte vi
som de egenlig konsekvente, finde vi fra Begyndelsen en
Stræben efter at hævde det Humanes og det Religiøses Eenhed
saaledes, at dette sidste sees som dybest begrundet i det menneskelige
Væsen og som bestemt i sin Form og sit Indhold ved dette.
Og det vil ogsaa være let at vise, at Erkjendelsen af Menneskets
Væsen i dets Fuldstændighed med Nødvendighed maa føre
til det religiøse Forhold, og at den kun kan afsluttes i dette.

Naar det for Menneskelivet Eiendommelige, det, der ved
en afgjørende Forskjel adskiller det fra de lavere Former af
Livet, fra Plantelivet og Dyrelivet, er Selvbevidsthedens
Bestemmelse, den frie Viden om sig selv og en Viden om et
Uendeligt, saa vises der herved tilbage til Menneskevæsenets
Uendelighed. Kun det Væsen, der har indre Uendelighed,
kan i sit Forhold til sig selv blive befriet fra den Bundethed
til det Enkelte, der karakteriserer Dyrets Selvfornemmelse,
som aldrig løsnes fra dets øieblikkelige Bestemthed; kun det
kan blive sig selv bevidst efter sin Eenhed og Almindelighed.
Og kun det Væsen, der har indre Uendelighed, kan have
Bevidsthed om et Uendeligt, thi denne forudsætter Bevidsthedens
Uendelighed, og denne atter det Væsens Uendelighed, hvis
Bevidsthed har denne Bestemmelse. Men det aandelige Væsen,
der har indre Uendelighed, kan kun naae sin Udfoldelses Maal,
kan kun finde Hvile og Tilfredsstillelse i et Uendelighedsforhold,
i et saadant Forhold, i hvilket dets Uendelighedsbestemmelse
virkeliggjøres, idet Væsenet udfolder sig som en aandelig
Livs-Eenhed, i hvilken enhver oprindelig Bestemthed, ethvert
virkeliggjort Anlæg, er gjort til Moment, draget ind under
% --- p. 499 ---
\opage{499}Væsenets bevidste Harmoni, der virkeliggjøres gjennem den
Aandsvirksomhed, som bliver det primitive Udtryk for det selvbevidste
Væsen: gjennem Erkjendelsen. Og dette Uendelighedsforhold,
der saaledes nærmest er bestemt som et Aandens Forhold
% emphasis: „til sig selv“ letterspaced (gaps 9--13 px; „kan“ after it 1--4 px)
\emph{til sig selv}, kan ikke gjennemtænkes, uden at det udfolder sig
til et Forhold \emph{til det Guddommelige}, til et \emph{religiøst}
Forhold. Naar vi søge den indre Uendeligheds \emph{Kilde} i
det Menneskelige, naar vi gjennemtænke Selvbevidsthedens
Begreb, komme vi med Nødvendighed til det Guddommelige.
Det menneskelige Selv, der bliver sig selv bevidst i sin indre
Uendelighed, kan kun blive sig bevidst som et \emph{afledet} Selv,
% p. 499: the print has „inde-|slutter“ (faint but clear e before the hyphen); the layer reads „ind-“
som et Selv, der har sin Grund i et Høiere, og saaledes indeslutter
selve Selvbevidstheden et Forhold til Menneskevæsenets
Kilde, til dets guddommelige Grund, der tillige er den hele Tilværelses
Grund. Den menneskelige Selvbevidsthed fuldbyrdes
først i Gudsbevidstheden. Fordi Menneskevæsenet grunder
i det Guddommelige, kan det som selvbevidst Væsen kun
komme til Virkeliggjørelse, idet dette Forhold, ved hvilket
Væsenet har Væren, bliver til for dets Bevidsthed og bliver
saaledes til for det, at Gudsforholdet gaaer i Eenhed med
Selvets Forhold til sig selv i Selvbevidstheden.

Saaledes bliver det religiøse Forholds indre Sammenhæng
med det Menneskeliges Væsensbestemmelse klar. Men
% p. 499: the last two letters of „Karakteer“ are very faintly printed at the line end; read as „er“ (cf. „Karakteer“ p. 495)
for at dette Forhold skal faae den \emph{sande} Uendeligheds Karakteer,
skal frigjøres fra den Endelighedsbegrændsning, der bestandig
bliver heftende ved de positive Religioner med deres fra Verden
sondrede, personliggjorte, undergjørende Guddom og med
deres Sondring af det religiøse Forhold fra Menneskeaandens
andre væsenlige Livsforhold, maa det blive et Forhold til et
Guddommeligt, der i Sandhed er et Uendeligt, idet det som
beherskende og omsluttende Tilværelsen er \emph{iboende} i den,
som begrundende Menneskeaanden er \emph{væsenslig} med den. Og
det maa blive et Forhold, der i Virkeligheden kan gjennemtrænge
\emph{alle} Aandens væsenlige Forhold, der kan finde sit
Udtryk i det \emph{sædelige} Livs Forhold, og som, medens det hæver
dem op i en høiere Sfære, selv vinder Fylde og Virkelighed
% --- p. 500 ---
\opage{500}gjennem dem. Saaledes kommer det religiøse Forhold til at
omfatte \emph{hele} Menneskelivet, \emph{alle} Aandens Livsyttringer; det
bliver det, i hvilket Aanden i hvert Øieblik finder Hvile og
Forsoning: Forsoning i sit eget Indre gjennem Væsenets
Harmoni, og Forsoning med sin Tilværelsesgrund.

Det saaledes bestemte religiøse Forhold, den \emph{ideale}
Religiøsitetens Form, den Religion, der kan blive \emph{Menneskehedens}
Religion, ligger altsaa i Menneskets Væsensbestemmelse
som en nødvendig, uafviselig Fordring; dette Forhold, taget
efter sin fulde, omfattende Betydning, bliver Menneskelivets
høieste, afsluttende Realisationsform.

Kun een Bestemmelse bliver endnu udtrykkelig at betone,
for at det her Sagte ikke skal misforstaaes. Religionens
Kilde er søgt i det menneskelige Væsen; den er seet som fremsprungen
af dette, og Religionen bliver fra dette Synspunkt seet
som Menneskeaandens Værk. Men idet det menneskelige Væsen
selv grunder i det Guddommelige, bliver Kildens Udspring
at søge dybere: i det Væsen, ved hvilket Menneskeaanden har
Uendeligheden som indre Væsensbestemmelse og Trangen til
Virkeliggjørelsen af denne Uendelighed. Fra dette Synspunkt
bliver da Religionen \emph{Guddommens Værk}, et Værk af
dens indre Aabenbaring for Menneskeaanden, d. v. s. af
dens \emph{Væsensmeddelelse} til den.

\begin{center}\rule{1.8cm}{0.4pt}\end{center}

Gjennem Opfattelsen af det \emph{religiøse Forhold} i dets
ideale, almeenmenneskelige Form er den Opfattelse af \emph{det
Guddommelige} antydet, til hvilken den humane Livsretning
konsekvent maa føre. Den kan ikke tilegne sig de
positive Religioners Forestilling om den personlige, undergjørende
Gud. I Personlighedens Bestemmelse, saaledes
som den gjøres gjældende i de positive Religioner og af
Theologien, kan den kun see en endeliggjørende Bestemmelse,
% printer's error: „Anthropomosisme“ (long s, no r; for „Anthropomorfisme“); transcribed as printed
en Anthropomosisme, der overføres paa Gud netop for at
holde ham i en Forestillingssondring fra det Menneskelige og
% --- p. 501 ---
\opage{501}fra Tilværelsen og for at frigjøre hans Handlen fra Fornuftnødvendighedens
Lov og skaffe Plads for Miraklet. Istedetfor
denne Bestemmelse af det Guddommelige maa den søge
en anden og høiere, der, medens den udtrykker Guds Væsen
% emphasis: the letterspacing runs through „Endeligheds“ into „in“ (gap after i 11 px at 300 dpi; before d borderline, 7 px) -- second reading; the rest of „-dskrænkningen“ set solid
som \emph{Aand}, frigjør det fra \emph{Endelighedsin}dskrænkningen;
der, medens den betegner Gud som Selvbevidsthedernes Kilde,
frigjør ham fra Selvbevidsthedsbestemmelsens Begrændsning
og relative Abstraktion. Og ved denne Bestemmelse vil den
da modsige Forestillingen om den \emph{undergjørende} Gud.
Den Gud, der ved Underet gjennembryder Naturens Love
og omstyrter Tankens, kan den tænkende Aand ikke anerkjende
uden at fornægte sit Væsen, som den ikke har fra sig selv,
men fra Guddommen. Den maa fordre en Gud, hvis Virken
ikke modsiger sig selv, hvad den undergjørende Guds Virken
% p. 501: a faint dot (risen space) after „sat,“; not transcribed
gjør, idet den ophæver det, den selv har sat, og saaledes viser
sig som endelig. Den maa fordre en Gud, hvis \emph{Væsen}
finder sit Udtryk i Fornuftens Love, der beherske Tilværelsen,
og som derfor ikke kan bryde disse Love.

Denne efter sit Væsens Fornuftnødvendighed handlende
Gud maa den saa, for at bevare hans Uendelighed, nødvendigviis
opfatte som \emph{iboende} i den Tilværelse, der har sin Grund
i Gud, og, som iboende i den, sammensluttende den i sin
Eenhed og saaledes \emph{ideelt} sondrende sig fra Tilværelsen,
saaledes som den er udfoldet i den \emph{reale} Værens Form.
Som iboende (immanent) i Tilværelsen, som den, der gjennemtrænger
Tilværelsen med sit Liv og indenfra behersker den,
maa Gud bestemmes som staaende i et \emph{Væsenseenheds} Forhold
til den reale Tilværelse og til Menneskeaanden. Fra Virkningens
Væsen maa der sluttes tilbage til Aarsagens, idet
Aarsagen kun kan være Aarsag til det, hvis Væsensbestemmelser
den bærer i sig. Som Ophav til Naturen maa Gud ikke
blot være et aandeligt, et ideelt Væsen, men have Realitetens
Bestemmelser i sig, være \emph{ideelt-realt} Væsen. Men som
Ophav til den Natur, der finder sin Afslutning i Liv og i
Aand, maa Gud være ideelt-realt Væsen i den Betydning,
at Idealiteten er den primitive Bestemmelse, Grundbestemmelsen
% --- p. 502 ---
\opage{502}i det Guddommelige, den, der omfatter og forklarer Realitetsbestemmelsen.
Ved Guds Væsen er da paa een Gang Naturen
i os hævdet i sin Væsenlighed og Aanden i sit ideale Herredømme.
Til eet Princip vise Tilværelsens Modsætninger
tilbage som til deres Kilde; i Naturens og i Aandens Verden
gjør overalt een guddommelig Fornuft sig gjældende, virkende
efter evige, ubrødelige Love, efter Love, der have deres Garanti
i selve Guddommens Væsen, hvis Udtryk de ere. Tilværelsens
Lovmæssighed er, som fremgaaet af den iboende evige Fornuft,
en evig gyldig Lovmæssighed. — Som Tilværelsens Maal
og som det, der bestemmer de endelige Aander og bliver deres
Handlings Maal og Norm, bliver endelig det Guddommelige
% emphasis: „ethiske“ letterspaced (gaps 8--12 px against 3--6 px solid)
at see under \emph{ethiske} Bestemmelser, som det Gode; men denne
ethiske Bestemmelse af det, der definitivt betegner det som
Aand, bliver paa det Nøieste forbunden med dets logiske og
fysiske Bestemmelser; det bliver kun disses fuldstændigere og
rigere Udtryk, ikke en Bestemmelse, der frigjør det Guddommelige
fra Tankens evige Love og, som Theologien vil,
skaffer en Plads for den undergjørende Vilkaarlighed.

% p. 502: a short centred rule (c. 1.7--1.8 cm) between the paragraphs; no section number
\begin{center}\rule{1.8cm}{0.4pt}\end{center}

Ved den Opfattelse af Menneskevæsenet og dets Forhold
til det Guddommelige, til hvilken den humane Livsretning
% emphasis: „Livsopgaven for Menne-|sket“ letterspaced across the line break (gaps 8--12 px; „ſket“ on the next line 10--14 px)
konsekvent maa føre, bestemmes \emph{Livsopgaven for Mennesket}.
Den kan kun være den frie og bevidste Virkeliggjørelse
af det menneskelige Væsen efter dets fuldstændige Bestemthed.
Men denne almene Formel for Livsopgaven maa faae et
fyldigere Udtryk. Vi kunne give den det, idet vi først
see det menneskelige Selv som det, der er en udelelig Eenhed
af Aand og Natur, og som det, i hvis forskjellige Grundvirksomheder
Væsenets Eenhed maa udtrykke sig i Samstemmen
med sig selv. Opgaven faaer da den bestemtere Form, at
Mennesket skal udvikle sin Personlighed i sine Grundvirksomheders
indre Harmoni og i uadskillelig Eenhed af det Aandelige
% emphasis: „Væsen“ letterspaced (gaps 8--10 px)
og det Naturlige. Mennesket faaer saaledes sit \emph{Væsen} til
% --- p. 503 ---
% emphasis: „konkrete“ letterspaced (gaps 7--11 px)
\opage{503}Opgave som den \emph{konkrete} Eenhed, i hvilken det reale
Naturgrundlag ikke skal fornægtes eller nedkues, men optages
i Aanden og gjennemformes af den, saaledes at Villien faaer
Indhold og Fylde fra sin Naturgrund, medens atter det,
der fremgaaer af Naturgrundlaget, faaer en Norm og en
Regel i den ethiske Fordring, der med det Oprindeliges (det
Aprioriskes) Præg fremgaaer af Aandens Væsen. Denne
ethiske Fordring, der har det Ubetingedes Præg, og som oprindelig
træder frem i abstrakt Form gjennem Bevidstheden
om en Forpligtelse overhovedet, om et ubetinget gyldigt:
„Du skal!“, skal saa, for at kunne gjennemforme Personlighedens
Naturgrundlag, gjennem Erkjendelsen af Menneskets
Væsen udfoldes i en Rigdom, i hvilken den svarer til Livets
og Handlingens Virkelighedsrigdom. Den oprindelige Bevidsthed
om det Gode, om Pligten, skal i Handlingens
Virkelighed blive til en Bevidsthed om dette bestemte Gode,
om denne bestemte Pligt, og under denne Udfoldelse skal den
bevare sin Vished, saa at Livets Heelhed heelt igjennem faaer
det klare og bevidste Ethiskes Præg.

Til det samme Maal vises vi ogsaa hen, naar vi see
den ovenfor givne almene Formel for den menneskelige Livsopgave
fra et andet Synspunkt. At Mennesket skal gjøre
% emphasis: „Væsen“ letterspaced
sit \emph{Væsen} til Opgave for sig efter dets alsidige og fuldstændige
Bestemthed, vil ikke blot sige, at det skal udvikle det
i harmonisk Eenhed af Aand og Natur og af Aandens Grundvirksomheder;
men det vil tillige sige, at det skal udvikle
% emphasis: „et eiendomme-|ligt Led i en Tilværelsens Heelhed“ letterspaced across the line break (gaps 10--14 px; „et“ gap 10 px against 5--6 px in „som“)
det, ikke som isoleret og atomistisk, men som \emph{et eiendommeligt
Led i en Tilværelsens Heelhed}, som indordnet i
denne Tilværelse og som underordnet under dens Princip.
Kun saaledes udvikler Selvet sig efter sin Fuldstændighed;
som løsrevet fra Livssammenhængen med Heelheden er det
netop det Ufuldstændige, ligesom det enkelte Organ i en
individuel Organisme mister sin Integritet, naar det skilles
fra den Heelhed, der bærer det. Men idet Selvet bliver sig
bevidst som indordnet i et mere Omfattende og som underordnet
% emphasis: „Hævden af sig selv“ letterspaced (gaps 11--15 px); the final one-letter „i“ at the line end cannot be judged and is left unmarked; „sin Eiendommelighed“ on p. 504 is solid
under et Høiere, bliver dets \emph{Hævden af sig selv} i
% --- p. 504 ---
% emphasis: „Hengivelse til det | Høiere“ letterspaced across the line break (gaps 9--13 px)
\opage{504}sin Eiendommelighed Et med dets \emph{Hengivelse til det
Høiere}, det skal tjene; med andre Ord: dets Selvhævdelse
% emphasis: „Ethiskes“ letterspaced (here and five lines below)
og Selvudfoldelse faaer det \emph{Ethiskes} Præg, dets Villen
af sig bliver en Villen af en høiere Fordring, af Tilværelsens
Lov og af Guddommens Villie. Saaledes taber Selvudviklingen
ethvert Præg af Egoisme; Udfoldelsen af den
individuelle Eiendommelighed bliver taget i Heelhedens og i
Principets Tjeneste; det Liv, Individet lever for sig, lever
det for den hele Menneskehed og for det Guddommelige.
Og idet Livet paa ethvert Punkt faaer det \emph{Ethiskes} Præg,
faaer det tillige, idet Selvet overalt i Tilværelsen, i alle de
% emphasis: „Religiøses“ letterspaced
relative Forhold, møder det Guddommelige, det \emph{Religiøses}
Præg; Livsheelheden bliver paa een Gang ethisk og religiøs;
det religiøse Forhold skiller sig ikke, som begrændset til enkelte
Tider eller til enkelte Gebeter af Livet, ud fra de andre
væsenlige Livsforhold, men omfatter, gjennemtrænger og
løfter alle de andre.

I denne Livsanskuelse, som hviler paa Overbeviisningen
om det Menneskeliges uendelige Værd, ifølge det Guddommeliges
Iboen i det i sin Grund uforstyrrede, universelle
Menneskevæsen, og som fører til Anerkjendelsen af det
Menneskelige i alle dets væsenlige Fremtrædelsesformer og
til en universel Kjærlighed til Menneskeheden, til det Menneskelige
% emphasis: „som“ and „Humane“ letterspaced
\emph{som} menneskeligt, er det \emph{Humane} saaledes i inderlig
Eenhed med det Religiøse, der sees som fremgaaende af
Menneskevæsenets Grund og som vindende Indhold igjennem
det Ethiske. Den viser energisk Mennesket hen til en Virksomhed
i Virkelighedens Verden, men til en Virksomhed, der
i Virkelighedens Gjerning bevarer Forbindelsen med det
Høiere. Ved i hvert Øieblik at minde om Solidariteten i
Menneskeslægten, minder den den Enkelte om Nødvendigheden
af, for at løse sin individuelle Opgave, at arbeide paa Løsningen
af de sociale Opgaver. Medens den Livsanskuelse,
der fornægter Naturen, kun samler Individerne i Arbeidet
for deres Salighed, men ikke kan forene dem til et Arbeide
for de jordiske Livsbetingelser for et Aandens Liv, samler
% --- p. 505 ---
\opage{505}den humane Livsanskuelse dem til dette Arbeide, og i selve
dette lader den dem arbeide for deres aandelige Maal. Kun
% emphasis: „en menneskeværdig | Tilværelse for Alle“ letterspaced across the line break
ved den kan Grundlaget lægges til \emph{en menneskeværdig
Tilværelse for Alle}, fordi kun den i det rette Forhold
anerkjender det Aandeliges og det Naturliges Betydning,
Individets Ret og Samfundets Fordring, Selvhensynet og
Solidariteten. —

Denne religiøs-humane Livsanskuelse danner sig dog ikke
et idyllisk Billede af Menneskelivet, et Utopien, frigjort fra
Virkelighedens Ufuldkommenheder. Den erkjender, ligeoverfor
det Sædeliges uendelige Fordring, Menneskelivets Afstand
fra dets Ideal, dets Ufuldkommenhed, dets Skyld, dets
Syndighed. Men idet den anerkjender den smertelige Afstand
mellem Ideal og Virkelighed, opgiver den ikke Troen paa
Mennesket, paa dets Evne til at virkeliggjøre det Gode, paa
% emphasis: „denne“ letterspaced
at give \emph{denne} Tilværelse et Evighedens Indhold. Naar
den positive Religion lærer Menneskevæsenets Forstyrrelse
ved Synden, en Forstyrrelse, som kun et guddommeligt Under
kan hæve, saa fastholder hiin Humanitetens Religiøsitet, at
% emphasis: only „-væsenet“ of „Menneskevæsenet“ is letterspaced (gaps 8--11 px; „Menneske“ solid, 1--6 px); marked as printed
Menneske\emph{væsenet} ikke kan være forvandlet ved en Handling,
der er fremgaaet af selve Væsenet, ligesaalidt som Organismens
Væsen forvandles, fordi dens Integritet er ophævet
under Sygdommen. Og ligesaalidt som den kan indrømme,
at Menneskevæsenet er forstyrret ved Synden, ligesaalidt kan
den indrømme, at Naturen ved Mennesket skulde være dragen
ind under Forstyrrelsen. Den kan det ikke, fordi den seer
Tilværelsens Heelhed som baaren af det iboende Princip.

Naar den anerkjender det Onde, ikke blot som en Ufuldkommenhed,
en Mangel, men som en ved Menneskets Villie
sat Modsætning til det Gode, saa seer den i den Villie, der
er virksom til det Onde, en Villie, der maa brydes i Kampen
mod det Gode, fordi den kæmper mod, hvad der er selve
% emphasis: „Væsenet“ letterspaced (gaps 9--12 px); „i det“ solid
\emph{Væsenet} i det Menneskelige, og maa fremkalde en Reaktion
fra dette. Og den seer tillige i denne Villie, der har den
falske Uendeligheds Skin, men som kun faaer dette Skin fra
det, den kæmper imod og stræber at fortrænge, et Virksomt,
% --- p. 506 ---
\opage{506}hvis Virkens Retning saaledes kan vendes om, at den i det
Ondes Tjeneste virkeliggjorte, uendelig anspændte Energi kan
blive virksom for det Gode, faae positivt Værd for dette. Ligesom
paa det Fysiskes Gebeet Kræfterne kunne omsættes i hverandre,
men i denne Omsætning Intet tabe, saaledes skal her
under Kraftens Forvandling Kraftsummen blive uforandret.
Men for at det Gode skal kunne virkeliggjøres i Individet, maa
ogsaa denne Livsanskuelse anerkjende, at Individerne trænge til
en Bistand. Denne Bistand søger den dog ikke udenfor den
Tilværelse, i hvilken Individet har sin Virkelighed; den søger
den i de andre, supplerende Individer, med hvilke Individet
slutter sig sammen i Kjærlighed, i gjensidig Væsensmeddelelse,
og, i sidste Instans, i det Princip, der er iboende i selve Individet,
og som slutter det sammen med de andre Individer
i Væsenets Eenhed. Idet Individerne træde i et Kjærlighedsforhold
til hverandre, seer den dem som indvirkende — paa
een Gang luttrende og supplerende — paa hverandre. Det
enkelte Individ bliver for det andet en Repræsentant for
Menneskets almene, fornuftbestemte Væsen, og som saadant
bliver det for Individet ligesom dets personificerede, objektiverede
Samvittighed. Idet det ene Individ i hiint Forhold
omsætter det andet til Repræsentant for det Idealt-Menneskelige,
faaer det i det en bestemmende Norm, en ideal Maalestok
for sin Handlen, et Korrektiv for den. Og idet Anskuelsen
af det Almeen-Menneskelige smelter sammen med
Anskuelsen af Personlighedernes Eiendommelighed, af det
Ethiskes eiendommelige Virkeliggjørelsesformer i de forskjellige
Individer, saa indeslutter den ved Kjærligheden, ved den
ubevidste Tro paa Mennesket idealiserede ethiske Eiendommelighed
ikke blot et Korrektiv for det i Individet Mangelfulde,
men tillige en Forvisning om den ethiske Kraft i
det Menneskelige, og de af Kjærlighedsforholdet omfattede
% printer's defect: „Saa|ledes“ broken at the line end with no hyphen printed; joined
Individer vise sig tillige som supplerende hverandre. Saaledes
bliver Kjærlighedsforholdet mellem Menneskene en ethisk
Livskilde for Individerne. Men gjennem dette Kjærlighedsforhold,
i hvilket Individerne sluttes sammen i en aandelig
% --- p. 507 ---
\opage{507}Eenhed, vises der hen til denne Eenheds Princip, til de individuelle
Aanders Udspring: det Guddommelige, og i dette
bliver i sidste Instans Bistanden at søge for Individet. Det
Guddommelige er, som den individuelle Aands Kilde, det,
der er det Gjennemvirkende i den; det er det, hvorfra
Menneskeaanden har den Uendelighedens Bestemmelse, der
træder frem i dens Forhold til sig selv, i dens Griben af
det Gode og selv i dens egoistiske Villen af sig selv i Modsætningen
til det Gode. Fra denne sin Grund kan Menneskeaanden
aldrig skilles, selv ikke naar den vil det, thi gjennem
selve den onde Villie virker den guddommelige Villie med
Væsenets Kraft som brydende og fortærende det Ondes Kraft.
I det Gjennembrud, i hvilket den gode Villie bryder frem
gjennem det Ondes Opløsning, er det Guddommelige virksomt;
thi Væsenet hævder sig kun i Gjennembruddet ved en
% emphasis: „det“ letterspaced (gaps 10, 8 px against 2--4 px in „Virken“)
Virken af \emph{det}, der er Væsenets Grund. Og i enhver Kamp,
i enhver Vaklen kan Villien søge tilbage til denne Grund,
der aabenbarer sig i Selvets Inderste gjennem den ethiske
Fordring, og den kan i den finde Fuldstændiggjørelsen af
den manglende Kraft gjennem den fulde Selvhengivelse, ved
hvilken det Eviges Kraft vindes, idet Individet frit taber
sig selv til det Evige. Her bliver ikke længer en Tvivl
% emphasis: „der“ letterspaced
om Forsoningen som \emph{der}, hvor Individet skjælvende veed sig
afhængigt af en ubegribelig Vilkaarligheds Valg eller Forkastelse,
% printer's error: no full stop after „(Calvin)“ before the next sentence „Her er Vished …“; transcribed as printed. „Calvin“ is in Fraktur, not antiqua
af Almagtens „forfærdelige Beslutning“ (Calvin)
Her er Vished om Forsoningen, fordi Individet finder det
Guddommelige i sig som iboende og bliver sig bevidst i sin
Væsenseenhed med det. Men idet saaledes den sidste Betingelse
for vort Livsmaals Virkeliggjørelse bliver det i os
virkende Guddommelige, vises vi atter tilbage til det Religiøse
som Fuldendelsen af vort ethiske Liv. I denne Eenhed af
det Ethiske og det Religiøse har hiint sin Sandhed, dette sin
Virkelighed.

% p. 507: a short centred rule (c. 1.8 cm) between the paragraphs; no section number
\begin{center}\rule{1.8cm}{0.4pt}\end{center}

Idet det saaledes er viist, at den humane Livsanskuelse
først kan afsluttes i det Religiøse, bliver derved dens normale
% --- p. 508 ---
\opage{508}Forhold til den positive Religion definitivt bestemt. Der maa
% emphasis: „Modsætning“ (first occurrence) and „den“ (two lines below) letterspaced
blive en \emph{Modsætning}, en klar og bevidst Modsætning;
men det i de positive Religioner Væsenlige kan trods \emph{den}
blive bevaret, og de positive, historiske Religioner kunne anerkjendes
% emphasis: „ideale“ and „Menneske-|hedens Religion“ letterspaced
i deres Betydning for Opnaaelsen af det \emph{ideale}
Religiøsitetens Trin, der er blevet betegnet som \emph{Menneskehedens
Religion}. Det er kun eensidige Former af den
humane Livsretning, der have kæmpet fjendtligt imod det
positivt Religiøse, og som have opfattet det som en blot
Forvanskning af Sandheden. Idet den humane Livsretning
vinder en dybere Erkjendelse af Menneskets Væsen og af
dets Grundforhold, kan den ikke undgaae at anerkjende, at
det netop er den samme Stræben efter Virkeliggjørelsen af
Menneskeaandens Væsen i et sandt Uendelighedsforhold, af
hvilken den ideale Religiøsitetens Form fremgaaer, som ogsaa
har været det dybest Bevægende og det Udfoldende i de positive
Religioner, om den end, idet den realiseredes under Endelighedsvilkaar,
i de positive Religioner endnu viser sig som
beheftet med en Endelighedsbegrændsning, der træder frem
i Opfattelsen af det Guddommelige, i det religiøse Forholds
Sondring fra de andre væsenlige menneskelige Forhold, i de
Former af de menneskelige Aandsvirksomheder, der komme til
Udfoldelse paa den positive Religions Trin, og i Forholdet
imellem dem. Overalt i de positive Religioner vil der derfor
kunne paavises ligesom en Brydning i Bevidstheden: Uendeligheden
vil momentviis gjennemtrænge Endeligheden, men
denne vil dog blive staaende som uophævet og hefte som et
Tyngende og Hæmmende ved Aandens Uendelighedsstræben.

Først i den Religiøsitetens Form, som udfolder sig af det
sig selv fuldt bevidste Humane, er Menneskeaandens oprindelige
Stræben frigjort fra den tyngende Endelighedsbegrændsning.
Det religiøse Forhold faaer den sande Uendelighedens Karakteer,
fordi det forbinder Menneskeaanden med et Guddommeligt, der
er iboende i det, og fordi det omfatter hele Menneskelivet, alle
de væsenlige menneskelige Forhold, saa at det Religiøse kan
blive Princip for det virkelige menneskelige Livs Totalitet.
% --- p. 509 ---
% p. 509: the first line is not indented; the paragraph continues from p. 508
\opage{509}I dette religiøse Forhold, og kun i det, kan den religiøse
Forsoning blive fuldstændig; thi kun i det har Aanden ikke
et guddommeligt Væsen overfor sig, der i sin ubegribelige
Villies Mysterium bliver det Fremmede for Aanden, men et
Guddommeligt, i hvilket Aanden gjenfinder sig selv og føler sig
bekræftet; — og kun i det har Aanden den fulde Forsoning
i sit eget Indre, idet alle Aandens væsenlige Virksomheder
ere i Harmoni, og idet Religionens Liv ikke leves ved en
Flugt fra Menneskelivets Virkelighed og ved en Fornægten
af væsenlige Yttringsformer af det. Som den Religiøsitet, der
gjennemtrænger hele Menneskelivet, kan denne Humanitetens
% emphasis: „Virkelighedens“, „indenfor Virkelig-|heden“ and „udeelte“ letterspaced
Religiøsitet betegnes som \emph{Virkelighedens} Religion og
som den, der bringer en Forsoning \emph{indenfor Virkeligheden},
en Forsoning for det \emph{udeelte} Menneskevæsen. Den
er Udtrykket for Religiøsitetens evige Aand som det hele
Menneskelivs Livsprincip. Den er som Religion i levende
indre Eenhed med Videnskab og Sædelighed. Den har sin
Rod i Menneskevæsenets dybeste Grund, og er derfor, idet
den frie og bevidste Aand har fundet sig selv og fundet
denne Religiøsitetens Form som Udtryk for sit Væsen, et
Urokkeligt og Uforgængeligt, af Tidernes Vexlen Uafhængigt.

Men netop Erkjendelsen heraf indeslutter Erkjendelsen
af den Sandhed, de positive Religioner indeslutte. Hvad der
er dybest begrundet i Menneskevæsenet, maa under Menneskelivets
hele historiske Udfoldelse have ladet sig tilsyne, om end
kun momentviis og ufuldstændigt. Aandens Uendelighedsstræben
maa derfor ogsaa aabenbare sig gjennem de positive
Religioner, og den vil gjøre det saameget desmere, som det
% emphasis: „mystiske“ letterspaced
\emph{mystiske} Element i Religionen træder frem, den af Forestillingen
bundne Reflexion træder tilbage. I det religiøse
Livs umiddelbareste Former, i den levende, barnlige, troende
Opfattelse kommer den religiøse Bevidsthed i sin dunkle
Trang Uendelighedsforholdet nærmest; hvor Reflexionen faaer
Magten, fjerner den sig derfra, og meest hvor Theologien,
den endelige Reflexion over Religionen, bliver Repræsentanten
for den religiøse Bevidsthed. Saa bliver Endeligheds%
% --- p. 510 ---
\opage{510}sondringen mellem den personlige Gud og Mennesket og
Naturen, Modsætningen mellem Guds Væsen og det Menneskelige,
Modsætningen mellem Guds undergjørende Virken
og Naturens og Tankens Love principmæssig gjennemført,
og med Endelighedsforholdet, der gjennem Sondringen og
Modsætningen fæstner sig i Bevidstheden, fæstnes Spliden,
% sic: „religtøse“ -- t for i (crossbar at x-height, no dot), p. 510; second reading, confirmed at 600 dpi
og Maalet, hvorefter den religtøse Bevidsthed stræber, bliver
ikke naaet. Bliver saa Theologien halvspekulativ — og
længer kan den ikke naae uden at ophøre at være Theologi —
saa opnaaes der kun ved den Mægling, som forsøges mellem
Tro og Viden, mellem det Antirationelle og det Rationelle,
mellem det positivt Religiøse og det Humane, at Bevidstheden
kan blive omtaaget og fordunklet, men der bliver ingen
Harmoni, ingen virkelig Eenhed. Saa maa den tænkende
Aand søge Eenheden med sig selv ad sin egen Vei, og den
finder den ved at fordybe sig i sit Væsen og ved, idet den
vinder det Humanes fuldstændige Begreb, at see det indeslutte
det Religiøse, der ikke strider mod Tanken, men bekræfter
Tanken, der ikke flygter fra Livet, men udfylder, forædler
og forhøier Livet. Her er saa Maalet fundet, og fra Maalet
kan der anerkjendende sees tilbage til de Stadier, der ligge
bag ved det; men Anerkjendelsen af deres relative Sandhed
er uadskillelig forenet med den klare og frie Bevidsthed om,
hvad der sondrer dem fra Maalet. Og derfor stræbes der
ikke heller efter at bevare Skinnet af den Eenhed, der i
Virkeligheden ikke er tilstede: fra den frie, humane Religiøsitets
Standpunkt maa ethvert nyrationalistisk Forsøg paa
efterat have afklædt Kristendommen Alt, hvad der gjør den
% emphasis: „det“ letterspaced
til positiv Religion, at hævde \emph{det} reducerede, abstrakt rationelle
Sublimat af Kristendommen, som man vedkjender sig,
det Kristeliges Navn, ubetinget forkastes som principløst og
forvirrende.

% p. 510: signature set right in bold Fraktur; then a centred double rule (c. 2.5 cm) closes the essay; nothing follows on the page
\begin{flushright}\textbf{Constans.}\end{flushright}
\begin{center}\rule{2.5cm}{0.4pt}\\[-10pt]\rule{2.5cm}{0.4pt}\end{center}

\end{document}
